\documentclass[11pt,a4paper]{article}

% --- Idioma y tipografía -------------------------------------------------
\usepackage[T1]{fontenc}
\usepackage[utf8]{inputenc}
\usepackage[spanish,es-noindentfirst,es-nodecimaldot,es-tabla]{babel}
\usepackage{lmodern}
\usepackage[scaled=0.95]{helvet}
\renewcommand{\familydefault}{\sfdefault}
\usepackage{microtype}
\usepackage{textcomp}
\usepackage{amssymb}
\usepackage{newunicodechar}
\newunicodechar{≥}{\ensuremath{\geq}}
\newunicodechar{≤}{\ensuremath{\leq}}
\newunicodechar{−}{\ensuremath{-}}
\newunicodechar{→}{\ensuremath{\rightarrow}}

% --- Página --------------------------------------------------------------
\usepackage[a4paper,top=2.6cm,bottom=2.4cm,left=2.3cm,right=2.3cm,headheight=26pt]{geometry}
\setlength{\parindent}{0pt}
\setlength{\parskip}{0.55em}
\linespread{1.08}
\setlength{\emergencystretch}{3em}

% --- Colores de marca ----------------------------------------------------
\usepackage[table]{xcolor}
\definecolor{sicharnaranja}{HTML}{F47920}
\definecolor{sicharnegro}{HTML}{1E1E1E}
\definecolor{sichargris}{HTML}{5A5A5A}
\definecolor{fondogris}{HTML}{F4F4F4}
\definecolor{rojo}{HTML}{C0392B}
\definecolor{ambar}{HTML}{D68910}
\definecolor{azul}{HTML}{2E5E8C}
\definecolor{verde}{HTML}{2E8B57}

% --- Paquetes generales --------------------------------------------------
\usepackage{graphicx}
\usepackage{longtable,booktabs,array,calc}
\usepackage{multirow}
\usepackage{enumitem}
\setlist{itemsep=0.15em,topsep=0.3em}
\setlist[itemize]{label=\textcolor{sicharnaranja}{\textbullet}}
\usepackage{etoolbox}
\usepackage{fancyhdr}
\usepackage{titlesec}
\usepackage{needspace}
\usepackage[most]{tcolorbox}
\usepackage[hidelinks,pdfusetitle]{hyperref}
\usepackage{bookmark}

% Soporte para las tablas que genera Pandoc
\providecommand{\tightlist}{\setlength{\itemsep}{0pt}\setlength{\parskip}{0pt}}
\newlength{\cslhangindent}
\makeatletter
\patchcmd\longtable{\par}{\if@noskipsec\mbox{}\fi\par}{}{}
\makeatother
\renewcommand{\arraystretch}{1.25}
\setlength{\LTpre}{0.6em}
\setlength{\LTpost}{0.8em}
\AtBeginEnvironment{longtable}{\small}
\arrayrulecolor{sichargris}
\setlength{\heavyrulewidth}{1pt}
\setlength{\lightrulewidth}{0.5pt}

% --- Títulos -------------------------------------------------------------
\setcounter{secnumdepth}{-\maxdimen}
\setcounter{tocdepth}{2}
\titleformat{\section}
  {\Large\bfseries\color{sicharnegro}}
  {}{0pt}
  {}[\vspace{-0.4em}{\color{sicharnaranja}\titlerule[1.2pt]}]
\titleformat{\subsection}{\large\bfseries\color{sicharnaranja}}{}{0pt}{}
\titleformat{\subsubsection}{\normalsize\bfseries\color{sicharnegro}}{}{0pt}{}
\titlespacing*{\section}{0pt}{1.6em}{0.9em}
\titlespacing*{\subsection}{0pt}{1.3em}{0.5em}
\titlespacing*{\subsubsection}{0pt}{1.1em}{0.3em}
\newcommand{\sectionbreak}{\clearpage}


% --- Encabezado y pie ----------------------------------------------------
\pagestyle{fancy}
\fancyhf{}
\fancyhead[L]{\small\color{sichargris}\textbf{SICHAR} · Consultoría estratégica completa}
\fancyhead[R]{\includegraphics[height=11pt]{logo-sichar.png}}
\fancyfoot[L]{\small\color{sichargris}Proyecto EA · Octubre de 2026}
\fancyfoot[R]{\small\color{sichargris}\thepage}
\renewcommand{\headrulewidth}{0.4pt}
\renewcommand{\headrule}{\hbox to\headwidth{\color{sicharnaranja}\leaders\hrule height \headrulewidth\hfill}}

% --- Cajas destacadas ----------------------------------------------------
\tcbset{
  cajabase/.style={enhanced, sharp corners, boxrule=0pt,
    leftrule=4pt, left=10pt, right=10pt, top=8pt, bottom=8pt,
    before skip=0.9em, after skip=0.9em,
    fonttitle=\bfseries, coltitle=sicharnegro,
    attach title to upper={\par\smallskip}},
}
\newtcolorbox{cajainfo}[1]{cajabase, colback=azul!6, colframe=azul, title={#1}}
\newtcolorbox{cajaaviso}[1]{cajabase, colback=ambar!9, colframe=ambar, title={#1}}
\newtcolorbox{cajaalerta}[1]{cajabase, colback=rojo!6, colframe=rojo, title={#1}}
\newtcolorbox{cajamarca}[1]{cajabase, colback=sicharnaranja!8, colframe=sicharnaranja, title={#1}}

% --- FODA ---------------------------------------------------------------
\tcbset{
  fodabase/.style={enhanced, sharp corners, boxrule=0.6pt, fonttitle=\bfseries\small,
    coltitle=white, left=6pt, right=6pt, top=4pt, bottom=4pt, fontupper=\small},
}
\newtcolorbox{fodaF}[1]{fodabase, colframe=verde, colback=verde!5, title={#1}}
\newtcolorbox{fodaD}[1]{fodabase, colframe=ambar, colback=ambar!6, title={#1}}
\newtcolorbox{fodaO}[1]{fodabase, colframe=azul, colback=azul!5, title={#1}}
\newtcolorbox{fodaA}[1]{fodabase, colframe=rojo, colback=rojo!5, title={#1}}

% --- Ficha de hallazgo ---------------------------------------------------
% #1 número, #2 nombre, #3 nombre sin formato (índice), #4 urgencia,
% #5 cuándo actuar, #6 responsable sugerido
\newcommand{\hallazgo}[6]{%
  \clearpage
  \phantomsection
  \addcontentsline{toc}{subsection}{Prioridad #1 · #3}%
  \begin{tcolorbox}[enhanced, sharp corners, boxrule=0pt, colback=sicharnegro,
      left=0pt, right=0pt, top=0pt, bottom=0pt, after skip=1em]
    \begin{tabular}{@{}>{\centering\arraybackslash}m{2.6cm}@{\hspace{10pt}}m{\linewidth-3.4cm}@{}}
      \cellcolor{sicharnaranja}\rule{0pt}{2.4cm}%
      \color{white}\begin{tabular}{@{}c@{}}\footnotesize\bfseries PRIORIDAD\\[2pt]
        \fontsize{38}{38}\selectfont\bfseries #1\end{tabular} &
      \color{white}{\small\itshape Hallazgo #1}\par\vspace{2pt}
      {\LARGE\bfseries #2}
    \end{tabular}
  \end{tcolorbox}
  \vspace{-0.6em}
  \noindent
  \begin{tcbraster}[raster columns=3, raster equal height, raster column skip=2mm]
    \begin{tcolorbox}[enhanced, sharp corners, boxrule=0pt, colback=rojo!10, fontupper=\small]
      {\scriptsize\bfseries\color{sichargris} URGENCIA}\par
      {\large\bfseries\color{rojo} #4}
    \end{tcolorbox}
    \begin{tcolorbox}[enhanced, sharp corners, boxrule=0pt, colback=fondogris, fontupper=\small]
      {\scriptsize\bfseries\color{sichargris} CUÁNDO ACTUAR}\par
      \textbf{#5}
    \end{tcolorbox}
    \begin{tcolorbox}[enhanced, sharp corners, boxrule=0pt, colback=fondogris, fontupper=\small]
      {\scriptsize\bfseries\color{sichargris} RESPONSABLE SUGERIDO}\par
      \textbf{#6}
    \end{tcolorbox}
  \end{tcbraster}
}

% --- Portada -------------------------------------------------------------
\newcommand{\portada}{%
  \begin{titlepage}
    \thispagestyle{empty}
    \begin{tikzpicture}[remember picture, overlay]
      \fill[sicharnegro] (current page.north west) rectangle ([yshift=-12cm]current page.north east);
      \fill[sicharnaranja] ([yshift=-12cm]current page.north west) rectangle ([yshift=-12.5cm]current page.north east);
    \end{tikzpicture}
    \vspace*{-1.2cm}
    \begin{tcolorbox}[enhanced, colback=white, boxrule=0pt, arc=3pt, width=5.2cm,
        left=6pt, right=6pt, top=6pt, bottom=6pt, halign=center]
      \includegraphics[height=1.4cm]{logo-sichar.png}\hspace{6pt}%
      \raisebox{0.45cm}{\Huge\bfseries\color{sicharnegro}SICHAR}
    \end{tcolorbox}
    \vspace{0.5cm}
    {\color{sicharnaranja}\small\bfseries PROYECTO EA · CONSULTORÍA ESTRATÉGICA}\par
    \vspace{0.3cm}
    {\color{white}\fontsize{30}{34}\selectfont\bfseries CONSULTORÍA\\ESTRATÉGICA\\COMPLETA\par}
    \vspace{2.6cm}
    {\fontsize{46}{50}\selectfont\bfseries\color{sicharnaranja} SICHAR\par}
    \vspace{0.6cm}
    {\Large\itshape\color{sichargris} Diagnóstico, priorización de hallazgos y hoja de ruta de
      soluciones para crecer sin perder la confianza del canal\par}
    \vfill
    \begin{tcolorbox}[enhanced, sharp corners, boxrule=0pt, leftrule=4pt,
        colframe=sicharnaranja, colback=fondogris, left=12pt, top=10pt, bottom=10pt]
      {\scriptsize\bfseries\color{sichargris} PREPARADO PARA}\par
      \textbf{Gerencia y socios de SICHAR} · Cali, Colombia\par
      \medskip
      {\scriptsize\bfseries\color{sichargris} PREPARADO POR}\par
      \textbf{Isabella Ruiz, Juan José Giraldo, Martin Osorio, David Arias, Mariana Restrepo y
        Karen Carvajal}\par
      \medskip
      \textit{Octubre de 2026}
    \end{tcolorbox}
  \end{titlepage}
  \renewcommand{\contentsname}{Contenido}
}

\hypersetup{
  pdftitle={Consultoría Estratégica Completa · SICHAR},
  pdfauthor={Isabella Ruiz, Juan José Giraldo, Martin Osorio, David Arias, Mariana Restrepo, Karen Carvajal},
  pdfsubject={Diagnóstico, priorización de hallazgos y hoja de ruta},
  pdflang={es-CO},
}

\begin{document}
\portada
\tableofcontents
\clearpage

\hypertarget{introducciuxf3n}{%
\section{1. Introducción}\label{introducciuxf3n}}

SICHAR es una empresa colombiana con más de 20 años de trayectoria en la importación y distribución
de acabados para la construcción. Actualmente, enfrenta desafíos como \textbf{el aumento del 30 \%
en los costos de importación desde India, el riesgo de conflictos con sus distribuidores, la falta
de segmentación de clientes, la baja madurez digital, la ausencia de asesoría especializada y la
poca presencia en redes sociales.}

Esta consultoría \textbf{busca diagnosticar y priorizar estas problemáticas, proponiendo soluciones
enfocadas en optimizar costos, fortalecer la relación con los distribuidores, definir perfiles de
clientes, mejorar los procesos comerciales mediante herramientas digitales, ofrecer asesoría
especializada y aumentar la visibilidad de la marca.}

Con la implementación de estas recomendaciones, SICHAR \textbf{podrá mejorar su eficiencia y
competitividad, fortalecer sus relaciones comerciales, aprovechar nuevas oportunidades de negocio y
aumentar su presencia en el mercado.} Para evaluar los avances, se plantean indicadores y metas a
doce meses que deberán validarse con los datos reales de la empresa.

\needspace{6\baselineskip}

\hypertarget{quuxe9-debe-esperar-usted-de-este-documento}{%
\subsection{1.1 Qué debe esperar usted de este
documento}\label{quuxe9-debe-esperar-usted-de-este-documento}}

Este documento está escrito para quien toma las decisiones, no para especialistas. Está pensado para
que usted pueda leerlo, decidir y actuar. Encontrará:

\begin{itemize}
\item
  \textbf{Un diagnóstico claro y honesto:} seis hallazgos, ordenados del más urgente al menos
  urgente, cada uno con su descripción, su impacto actual y su impacto potencial.
\item
  \textbf{Ejemplos reales de empresas que ya resolvieron problemas parecidos} (nacionales e
  internacionales) y qué podría tomar SICHAR de cada una.
\item
  \textbf{Soluciones tangibles:} acciones concretas, responsables sugeridos, plazos, entregables e
  indicadores para cada hallazgo.
\item
  \textbf{Una hoja de ruta integrada} a 30 días, 90 días, 6 y 12 meses, y la lista de decisiones que
  se necesitan de la gerencia.
\end{itemize}

\emph{Lo que este documento no es:} no es un plan financiero detallado ni una auditoría contable.
Las metas que se proponen son \textbf{sugeridas y deben validarse con la línea base real} de SICHAR
durante el primer mes de ejecución.

\needspace{12\baselineskip}

\hypertarget{alcance-del-proyecto}{%
\subsection{1.2 Alcance del proyecto}\label{alcance-del-proyecto}}

\begin{longtable}[]{@{}
  >{\raggedright\arraybackslash}p{(\columnwidth - 2\tabcolsep) * \real{0.5000}}
  >{\raggedright\arraybackslash}p{(\columnwidth - 2\tabcolsep) * \real{0.5000}}@{}}
\toprule\noalign{}
\begin{minipage}[b]{\linewidth}\raggedright
\textbf{Incluye}
\end{minipage} & \begin{minipage}[b]{\linewidth}\raggedright
\textbf{No incluye}
\end{minipage} \\
\midrule\noalign{}
\endfirsthead
\bottomrule\noalign{}
\endlastfoot
\begin{minipage}[t]{\linewidth}\raggedright
\begin{itemize}
\item
  Diagnóstico estratégico de SICHAR: fortalezas, oportunidades, debilidades y amenazas.
\item
  Identificación y priorización de seis hallazgos negativos.
\item
  Benchmarking con empresas de referencia para cada hallazgo.
\item
  Propuesta de soluciones accionables y de indicadores de seguimiento.
\item
  Hoja de ruta de 12 meses.
\end{itemize}
\end{minipage} & \begin{minipage}[t]{\linewidth}\raggedright
\begin{itemize}
\item
  Estados financieros, proyecciones de ventas o valoración de la empresa.
\item
  Implementación operativa de las soluciones (salvo acuerdo posterior).
\item
  Selección o compra de software específico.
\item
  Asesoría legal o tributaria sobre importaciones.
\end{itemize}
\end{minipage} \\
\end{longtable}

\textbf{Cobertura:} las cuatro líneas de producto de SICHAR (accesorios de baño, cerámicas y
acabados, elementos decorativos y remodelación/adecuación con referencias importadas), sus cinco
sedes (Cali, Medellín, Bogotá, Barranquilla y Bucaramanga) su red de proveedores en cinco países
(China, India, Brasil, España y Perú) y su canal de distribuidores y ferreterías. Los segmentos
nuevos analizados son \textbf{constructoras, arquitectos y empresas de proyectos}.

\needspace{6\baselineskip}

\hypertarget{lo-que-sichar-ya-tiene-a-su-favor}{%
\subsection{1.3 Lo que SICHAR ya tiene a su favor}\label{lo-que-sichar-ya-tiene-a-su-favor}}

Antes de hablar de problemas, conviene reconocer los cimientos sobre los que se construirán las
soluciones:

\begin{itemize}
\item
  Trayectoria de \textbf{más de 20 años} y un modelo B2B probado con distribuidores y ferreterías en
  todo el país.
\item
  Red de proveedores \textbf{directa en 5 países}, sin intermediarios de importación.
\item
  \textbf{Cobertura nacional} con 5 sedes estratégicas y equipo comercial propio en cada una.
\item
  Un \textbf{producto insignia bien posicionado}: el porcelanato brillante 60x130.
\item
  Reconocimiento de marca sostenido con voz a voz, \textbf{sin inversión en publicidad}.
\item
  \textbf{Identificación proactiva} de nuevos segmentos: constructoras, arquitectos y empresas de
  proyectos.
\end{itemize}

\hypertarget{lo-muxe1s-grave-que-encontramos}{%
\section{2. Lo más grave que encontramos}\label{lo-muxe1s-grave-que-encontramos}}

\begin{cajaalerta}{El hallazgo más grave: el riesgo de conflicto de canal}

La idea de \textbf{vender directo al consumidor final} para aumentar márgenes pone en riesgo a los
clientes que \textbf{hoy sostienen la operación}: ferreterías y distribuidores. Si ellos perciben
que SICHAR su proveedor les quita clientela, pueden reducir pedidos, priorizar a la competencia y
propagar su inconformidad por voz a voz, el único motor de reputación que SICHAR ha tenido durante
20 años. \textbf{Vender a constructoras y arquitectos también es conflicto de canal.} Si SICHAR les
factura directo, compite con los distribuidores que quizá hoy atienden esas obras. Falta una regla:
a partir de qué tamaño de proyecto la cuenta es directa, cómo se registra el proyecto y qué comisión
recibe el distribuidor de la zona.

\end{cajaalerta}

\textbf{¿Por qué es el más grave?} Porque es el único problema que, mal manejado, puede destruir el
\textbf{activo principal} de la empresa (la confianza del canal) y porque es una decisión que
\textbf{todavía está a tiempo de tomarse bien}. Los otros cinco hallazgos son debilidades que se
corrigen con trabajo; este es un riesgo que, si se comete el error, puede ser irreversible.

\textbf{En segundo lugar} está la exposición a los costos de importación: el alza del \textbf{30 \%
en los productos de India} ya presiona el margen y, si no se actúa, hace inviable financiar el resto
de las mejoras. Estos dos hallazgos juntos ---uno sobre el canal, otro sobre el margen son el eje
del plan.

\needspace{6\baselineskip}

\subsection{2.1 Cómo priorizamos}

Calificamos cada hallazgo de 1 a 5 en cuatro criterios: impacto en ventas o margen (35 \%), urgencia
(30 \%), cuántas otras soluciones dependen de él (20 \%) y facilidad para empezar con bajo costo (15
\%).

\begin{longtable}[]{@{}
  >{\raggedright\arraybackslash}p{(\columnwidth - 10\tabcolsep) * \real{0.2458}}
  >{\raggedright\arraybackslash}p{(\columnwidth - 10\tabcolsep) * \real{0.1421}}
  >{\raggedright\arraybackslash}p{(\columnwidth - 10\tabcolsep) * \real{0.1545}}
  >{\raggedright\arraybackslash}p{(\columnwidth - 10\tabcolsep) * \real{0.1628}}
  >{\raggedright\arraybackslash}p{(\columnwidth - 10\tabcolsep) * \real{0.1587}}
  >{\raggedright\arraybackslash}p{(\columnwidth - 10\tabcolsep) * \real{0.1360}}@{}}
\toprule\noalign{}
\begin{minipage}[b]{\linewidth}\raggedright
\textbf{Hallazgo}
\end{minipage} & \begin{minipage}[b]{\linewidth}\raggedright
\textbf{Impacto}
\end{minipage} & \begin{minipage}[b]{\linewidth}\raggedright
\textbf{Urgencia}
\end{minipage} & \begin{minipage}[b]{\linewidth}\raggedright
\textbf{Habilita a otros}
\end{minipage} & \begin{minipage}[b]{\linewidth}\raggedright
\textbf{Facilidad}
\end{minipage} & \begin{minipage}[b]{\linewidth}\raggedright
\textbf{Puntaje}
\end{minipage} \\
\midrule\noalign{}
\endfirsthead
\bottomrule\noalign{}
\endlastfoot
1. Conflicto de canal & 5 & 5 & 5 & 4 & \textbf{4,85} \\
2. Costos de importación & 5 & 5 & 3 & 3 & \textbf{4,30} \\
3. Buyer persona & 3 & 3 & 5 & 5 & \textbf{3,70} \\
4. Madurez digital & 4 & 3 & 4 & 2 & \textbf{3,40} \\
5. Atención especializada & 4 & 2 & 2 & 3 & \textbf{2,85} \\
6. Redes sociales & 3 & 2 & 1 & 4 & \textbf{2,45} \\
\end{longtable}

\needspace{12\baselineskip}

\hypertarget{mapa-de-prioridades}{%
\subsubsection{Mapa de prioridades}\label{mapa-de-prioridades}}

\begin{longtable}[]{@{}
  >{\raggedright\arraybackslash}p{(\columnwidth - 6\tabcolsep) * \real{0.1091}}
  >{\raggedright\arraybackslash}p{(\columnwidth - 6\tabcolsep) * \real{0.3175}}
  >{\raggedright\arraybackslash}p{(\columnwidth - 6\tabcolsep) * \real{0.1290}}
  >{\raggedright\arraybackslash}p{(\columnwidth - 6\tabcolsep) * \real{0.4444}}@{}}
\toprule\noalign{}
\begin{minipage}[b]{\linewidth}\raggedright
\textbf{Prioridad}
\end{minipage} & \begin{minipage}[b]{\linewidth}\raggedright
\textbf{Hallazgo}
\end{minipage} & \begin{minipage}[b]{\linewidth}\raggedright
\textbf{Urgencia}
\end{minipage} & \begin{minipage}[b]{\linewidth}\raggedright
\textbf{Por qué en este orden}
\end{minipage} \\
\midrule\noalign{}
\endfirsthead
\bottomrule\noalign{}
\endlastfoot
\textbf{1} & Riesgo de conflicto de canal & \textbf{Crítica} & Amenaza al canal que sostiene el
negocio; condiciona la forma de ejecutar casi todo lo demás. \\
\textbf{2} & Exposición a alzas en costos de importación & \textbf{Crítica} & Efecto ya en marcha
sobre el margen y la caja; financia las demás mejoras. \\
\textbf{3} & Buyer persona no definido & \textbf{Alta} & Cimiento de las acciones comerciales y de
marketing; barato y rápido. \\
\textbf{4} & Baja madurez tecnológica y digital & \textbf{Alta} & Infraestructura para la Vitrina,
el catálogo y las redes. \\
\textbf{5} & Falta de atención especializada & \textbf{Media-alta} & Diferenciador para arquitectos
y constructoras; requiere los perfiles y el catálogo. \\
\textbf{6} & Baja presencia en redes sociales & \textbf{Media} & Alta visibilidad, pero depende de
tener definidos canal, clientes y catálogo. \\
\end{longtable}

\hypertarget{quuxe9-hay-que-modificar}{%
\section{3. Qué hay que modificar}\label{quuxe9-hay-que-modificar}}

Resumen en una página de \textbf{los cambios concretos} que proponemos, uno por hallazgo. El detalle
de cada uno está en la sección 6.

\begin{longtable}[]{@{}
  >{\raggedright\arraybackslash}p{(\columnwidth - 6\tabcolsep) * \real{0.0496}}
  >{\raggedright\arraybackslash}p{(\columnwidth - 6\tabcolsep) * \real{0.2282}}
  >{\raggedright\arraybackslash}p{(\columnwidth - 6\tabcolsep) * \real{0.4345}}
  >{\raggedright\arraybackslash}p{(\columnwidth - 6\tabcolsep) * \real{0.2877}}@{}}
\toprule\noalign{}
\begin{minipage}[b]{\linewidth}\raggedright
\textbf{\#}
\end{minipage} & \begin{minipage}[b]{\linewidth}\raggedright
\textbf{Hoy}
\end{minipage} & \begin{minipage}[b]{\linewidth}\raggedright
\textbf{Qué hay que modificar}
\end{minipage} & \begin{minipage}[b]{\linewidth}\raggedright
\textbf{Primer paso (próximos 30 días)}
\end{minipage} \\
\midrule\noalign{}
\endfirsthead
\bottomrule\noalign{}
\endlastfoot
\textbf{1} & Se quiere abrir el B2C sin reglas con el canal. & Adoptar el modelo \textbf{``Vitrina
SICHAR + Distribuidor Aliado''}: la venta al consumidor final pasa por los distribuidores; la venta
directa solo con productos y proyectos que no compiten con ellos. & Mapear el canal y entrevistar a
8--10 distribuidores clave. \\
\textbf{2} & Costos de India +30 \% sin análisis por referencia. & Calcular el \textbf{costo total
puesto en bodega por referencia}, diversificar origen y fijar una \textbf{política de precios por
escalones}. & Construir la hoja de costos de las 20 referencias top. \\
\textbf{3} & Clientes conocidos ``de memoria''. & Definir \textbf{3--4 buyer personas} validadas con
entrevistas y registrar cada visita en un formulario único. & Agendar 15--20 entrevistas de clientes
y prospectos. \\
\textbf{4} & Procesos manuales; sin estrategia digital. & Ejecutar la \textbf{Ruta Digital} en 3
etapas: ordenar datos, abrir la puerta (web y catálogo), autoservicio. & Designar un líder digital y
mapear cuellos de botella. \\
\textbf{5} & Se vende producto sin asesoría. & Crear la \textbf{Mesa de Asesoría SICHAR} y el
programa ``Arquitecto Aliado'', con piloto en Cali. & Definir el perfil del asesor y diseñar la guía
de acabados. \\
\textbf{6} & Presencia mínima en redes. & Activar un \textbf{plan de contenido de 90 días} (3
publicaciones por semana) centrado en proyectos reales. & Auditar la cuenta actual y la de Cerámica
Italia. \\
\end{longtable}

\hypertarget{quuxe9-pasaruxe1-despuuxe9s-de-aplicar-las-recomendaciones}{%
\section{4. Qué pasará después de aplicar las
recomendaciones}\label{quuxe9-pasaruxe1-despuuxe9s-de-aplicar-las-recomendaciones}}

Si SICHAR ejecuta la hoja de ruta, la empresa pasa de ser \textbf{un proveedor tradicional que crece
por voz a voz} a ser \textbf{el proveedor aliado, visible y asesor} de ferreterías, distribuidores,
constructoras y arquitectos. Así se vería el cambio:

\begin{longtable}[]{@{}
  >{\raggedright\arraybackslash}p{(\columnwidth - 4\tabcolsep) * \real{0.2183}}
  >{\raggedright\arraybackslash}p{(\columnwidth - 4\tabcolsep) * \real{0.2877}}
  >{\raggedright\arraybackslash}p{(\columnwidth - 4\tabcolsep) * \real{0.4940}}@{}}
\toprule\noalign{}
\begin{minipage}[b]{\linewidth}\raggedright
\textbf{Dimensión}
\end{minipage} & \begin{minipage}[b]{\linewidth}\raggedright
\textbf{Hoy}
\end{minipage} & \begin{minipage}[b]{\linewidth}\raggedright
\textbf{Después de aplicar las recomendaciones}
\end{minipage} \\
\midrule\noalign{}
\endfirsthead
\bottomrule\noalign{}
\endlastfoot
\textbf{Relación con el canal} & Riesgo de ruptura si se abre el B2C sin reglas. & Distribuidores
que \textbf{ganan} con la estrategia digital: reciben referidos, muestrarios, capacitación y
protección de proyectos. El B2C se convierte en una fuente de ventas \textbf{para el canal}, no
contra él. \\
\textbf{Margen} & Presionado por el alza en India. & Costo total conocido por referencia, menos
dependencia de un solo origen y precios ajustados con método y por escalones. \\
\textbf{Conocimiento del cliente} & Disperso en la memoria de los vendedores. & 3--4 buyer personas
validadas y una \textbf{base de datos viva} de clientes y prospectos. \\
\textbf{Operación digital} & Procesos manuales. & Catálogo digital, cotizaciones en menos de 48
horas, CRM en uso en las 5 sedes y pedidos de distribuidores por canal digital. \\
\textbf{Diferenciación} & Se compite por producto y precio. & SICHAR ofrece \textbf{asesoría de
diseño y muestras} a arquitectos y constructoras: compite con acompañamiento. \\
\textbf{Visibilidad} & Reconocimiento solo por voz a voz. & Presencia constante en Instagram,
Pinterest, LinkedIn y WhatsApp, con proyectos reales y un ``Dónde comprar'' que alimenta al
canal. \\
\textbf{Nuevos segmentos} & Identificados pero sin estrategia. & Primeras cuentas de
\textbf{constructoras, arquitectos y firmas de proyectos} con ofertas propias para cada perfil. \\
\end{longtable}

\needspace{12\baselineskip}

\hypertarget{resultados-esperados-a-12-meses-metas-sugeridas-por-validar-con-la-luxednea-base}{%
\subsubsection{Resultados esperados a 12 meses (metas sugeridas, por validar con la línea
base)}\label{resultados-esperados-a-12-meses-metas-sugeridas-por-validar-con-la-luxednea-base}}

\begin{longtable}[]{@{}
  >{\raggedright\arraybackslash}p{(\columnwidth - 2\tabcolsep) * \real{0.6448}}
  >{\raggedright\arraybackslash}p{(\columnwidth - 2\tabcolsep) * \real{0.3552}}@{}}
\toprule\noalign{}
\begin{minipage}[b]{\linewidth}\raggedright
\textbf{Resultado}
\end{minipage} & \begin{minipage}[b]{\linewidth}\raggedright
\textbf{Meta sugerida}
\end{minipage} \\
\midrule\noalign{}
\endfirsthead
\bottomrule\noalign{}
\endlastfoot
Retención de los 20 distribuidores principales & ≥ 95 \% \\
Quejas críticas por conflicto de canal & 0 \\
Referencias críticas con 2 o más proveedores validados & 100 \% \\
Tiempo de respuesta a cotizaciones & \textless{} 48 horas (\textless{} 24 horas en referencias
top) \\
Visitas de ventas registradas en el CRM & ≥ 80 \% \\
Proyectos asesorados por la Mesa de Asesoría (piloto de 90 días) & 10--15 \\
Publicaciones en redes & 3 por semana, de forma sostenida \\
\end{longtable}

\begin{cajaaviso}{Importante}

Estas metas son \textbf{propuestas de partida}, no promesas. SICHAR debe medir en el primer mes su
situación real (ventas por cliente, margen por línea, tiempos de cotización, etc.) y ajustar cada
meta con esa línea base.

\end{cajaaviso}

\hypertarget{anuxe1lisis-foda}{%
\section{5. Análisis FODA}\label{anuxe1lisis-foda}}

El FODA resume la situación de SICHAR hoy: lo que tiene a favor (\textbf{Fortalezas} y
\textbf{Oportunidades}) y lo que debe corregir o vigilar (\textbf{Debilidades} y \textbf{Amenazas}).

\begin{tcbraster}[raster columns=2, raster equal height=rows, raster column skip=3mm, raster row skip=3mm]
\begin{fodaF}{FORTALEZAS (internas, a favor)}
\begin{itemize}
\item
  Más de \textbf{20 años} de trayectoria y experiencia en el sector
  construcción.
\item
  Relación \textbf{directa con proveedores internacionales} en 5 países,
  sin intermediarios.
\item
  \textbf{Cobertura nacional} consolidada: 5 sedes estratégicas con
  equipo propio.
\item
  \textbf{Producto insignia} reconocido: porcelanato brillante 60x130.
\end{itemize}
\end{fodaF}
\begin{fodaD}{DEBILIDADES (internas, a corregir)}
\begin{itemize}
\item
  Ausencia de un \textbf{buyer persona} y de segmentación clara.
\item
  \textbf{Baja madurez digital:} sin estrategia de marketing ni
  presencia relevante en redes.
\item
  Falta de \textbf{atención especializada} (diseño) para nuevos
  segmentos.
\item
  Fuerza comercial basada en un \textbf{modelo tradicional de
  correrías}, con potencial limitado de escalar.
\end{itemize}
\end{fodaD}
\begin{fodaO}{OPORTUNIDADES (externas, a aprovechar)}
\begin{itemize}
\item
  \textbf{Segmentos sin explotar:} constructoras, arquitectos y empresas
  de proyectos.
\item
  \textbf{Canales digitales y comercio electrónico} para llegar al
  consumidor final.
\item
  Posibilidad de ofrecer \textbf{asesoría especializada en diseño} como
  diferenciador frente a la competencia.
\end{itemize}
\end{fodaO}
\begin{fodaA}{AMENAZAS (externas, a vigilar)}
\begin{itemize}
\item
  \textbf{Aumento de costos de importación} (India +30 \%) que presiona
  márgenes.
\item
  Riesgo de que los \textbf{proveedores no mantengan} las condiciones de
  negociación originales.
\item
  \textbf{Conflicto de canal} con los distribuidores actuales al avanzar
  hacia B2C.
\item
  Competencia directa y establecida de \textbf{Cerámica Italia}.
\end{itemize}
\end{fodaA}
\end{tcbraster}

\needspace{6\baselineskip}

\hypertarget{estrategias-cruzadas-cuxf3mo-usar-el-foda-para-decidir}{%
\subsection{5.1 Estrategias cruzadas: cómo usar el FODA para
decidir}\label{estrategias-cruzadas-cuxf3mo-usar-el-foda-para-decidir}}

Cruzando los cuatro cuadrantes obtenemos las estrategias que sustentan las soluciones de este
documento:

\begin{longtable}[]{@{}
  >{\raggedright\arraybackslash}p{(\columnwidth - 4\tabcolsep) * \real{0.1885}}
  >{\raggedright\arraybackslash}p{(\columnwidth - 4\tabcolsep) * \real{0.5060}}
  >{\raggedright\arraybackslash}p{(\columnwidth - 4\tabcolsep) * \real{0.3056}}@{}}
\toprule\noalign{}
\begin{minipage}[b]{\linewidth}\raggedright
\textbf{Cruce}
\end{minipage} & \begin{minipage}[b]{\linewidth}\raggedright
\textbf{Estrategia}
\end{minipage} & \begin{minipage}[b]{\linewidth}\raggedright
\textbf{Se concreta en\ldots{}}
\end{minipage} \\
\midrule\noalign{}
\endfirsthead
\bottomrule\noalign{}
\endlastfoot
\textbf{FO}

Fortalezas + Oportunidades & \begin{minipage}[t]{\linewidth}\raggedright
\begin{itemize}
\item
  Usar la red de 5 sedes y el producto insignia para llegar a arquitectos y constructoras.
\item
  Aprovechar los proveedores directos para ofrecer colecciones de proyecto.
\end{itemize}
\end{minipage} & Hallazgos 3 y 5: buyer personas, Mesa de Asesoría y programa ``Arquitecto
Aliado'' \\
\textbf{DO}

Debilidades + Oportunidades & \begin{minipage}[t]{\linewidth}\raggedright
\begin{itemize}
\item
  Cerrar la brecha digital para aprovechar el comercio electrónico y las redes.
\item
  Convertir la falta de asesoría en un servicio diferenciador.
\end{itemize}
\end{minipage} & Hallazgos 4, 5 y 6: Ruta Digital, Mesa de Asesoría y plan de contenido \\
\textbf{FA}

Fortalezas + Amenazas & \begin{minipage}[t]{\linewidth}\raggedright
\begin{itemize}
\item
  Usar los 5 países de proveedores para diversificar el origen frente al alza de India.
\item
  Usar la confianza de 20 años con los distribuidores para pactar reglas de canal antes del B2C.
\end{itemize}
\end{minipage} & Hallazgos 1 y 2: política de canal y diversificación de proveedores \\
\textbf{DA}

Debilidades + Amenazas & \begin{minipage}[t]{\linewidth}\raggedright
\begin{itemize}
\item
  Convertir el conocimiento disperso de los clientes en datos para anticipar cambios de precio y de
  mercado.
\item
  Fortalecer la presencia digital para no perder terreno frente a Cerámica Italia.
\end{itemize}
\end{minipage} & Hallazgos 2, 3 y 6: costos por referencia, CRM y redes \\
\end{longtable}

\hypertarget{hallazgos-en-orden-de-prioridad}{%
\section{6. Hallazgos en orden de prioridad}\label{hallazgos-en-orden-de-prioridad}}

A continuación, cada hallazgo se presenta con la misma estructura, de modo que pueda comparar y
decidir con facilidad:

\begin{itemize}
\item
  \textbf{1. Descripción detallada} del hallazgo.
\item
  \textbf{2. Impacto} que tiene hoy y que podría tener.
\item
  \textbf{3. Modelos inspiracionales:} empresas que ya resolvieron un problema similar.
\item
  \textbf{4. Solución tangible} para SICHAR: acciones, responsables, plazos y entregables.
\item
  \textbf{5. Indicadores} para medir si funciona.
\end{itemize}

\begin{cajainfo}{Cómo leer los responsables y plazos}

Los \textbf{responsables} son \textbf{sugeridos} según los cargos típicos de una empresa como
SICHAR; la gerencia puede reasignarlos. Los \textbf{plazos} se cuentan desde el inicio del proyecto
de implementación (semana 1).

\end{cajainfo}

\hallazgo{1}{\textbf{Riesgo de conflicto de canal}}{Riesgo de conflicto de canal}{CRÍTICA}{\textbf{Antes de lanzar cualquier venta directa (0--30 días)}}{\textbf{Gerencia general + Dirección comercial}}

\needspace{6\baselineskip}

\hypertarget{descripciuxf3n-detallada-del-hallazgo}{%
\subsection*{1. Descripción detallada del hallazgo}\label{descripciuxf3n-detallada-del-hallazgo}}

SICHAR ha crecido más de 20 años con una cadena simple y probada: \textbf{SICHAR → distribuidores y
ferreterías → consumidor final}. Su propósito de marca lo resume así: \emph{``Somos su proveedor, no
su competidor''}. Hoy esos distribuidores y ferreterías son \textbf{los clientes que sostienen la
operación}.

El empresario quiere incorporar un modelo \textbf{B2C} (venta directa al consumidor final) para
capturar un mayor margen. El problema es que, si SICHAR vende al mismo consumidor que hoy atienden
sus clientes, \textbf{pasa de ser su proveedor a ser su competidor}: los distribuidores verían cómo
la empresa que les surte les quita clientela, y con razón lo interpretarían como un conflicto de
intereses.

Este es un riesgo de \textbf{doble filo}: el B2C promete más margen, pero el canal que se podría
perder es el que genera hoy la mayor parte de los ingresos. Además, como la reputación de SICHAR se
ha construido casi exclusivamente por \textbf{voz a voz}, una molestia del canal se propagaría
rápido entre ferreterías y distribuidores de las cinco sedes.

\needspace{12\baselineskip}

\hypertarget{impacto-que-tiene-y-que-podruxeda-tener}{%
\subsection*{2. Impacto que tiene y que podría
tener}\label{impacto-que-tiene-y-que-podruxeda-tener}}

\begin{longtable}[]{@{}
  >{\raggedright\arraybackslash}p{(\columnwidth - 2\tabcolsep) * \real{0.5000}}
  >{\raggedright\arraybackslash}p{(\columnwidth - 2\tabcolsep) * \real{0.5000}}@{}}
\toprule\noalign{}
\begin{minipage}[b]{\linewidth}\raggedright
\textbf{Impacto actual (lo que ya está pasando)}
\end{minipage} & \begin{minipage}[b]{\linewidth}\raggedright
\textbf{Impacto potencial (si no se actúa)}
\end{minipage} \\
\midrule\noalign{}
\endfirsthead
\bottomrule\noalign{}
\endlastfoot
\begin{minipage}[t]{\linewidth}\raggedright
\begin{itemize}
\item
  Hay una decisión estratégica pendiente (abrir o no el B2C) y ninguna regla escrita sobre cómo
  convivirían ambos canales.
\item
  La incertidumbre ya condiciona la planeación comercial y la comunicación con los distribuidores.
\end{itemize}
\end{minipage} & \begin{minipage}[t]{\linewidth}\raggedright
\begin{itemize}
\item
  \textbf{Pérdida de ventas} del canal que sostiene el negocio: menos pedidos, menor prioridad a las
  referencias SICHAR o cambio hacia proveedores competidores como Cerámica Italia.
\item
  \textbf{Ruptura de la promesa de marca} (``proveedor, no competidor'') y daño reputacional que se
  propaga por voz a voz.
\item
  Guerra de precios entre SICHAR y sus propios clientes, con erosión del margen de todos.
\end{itemize}
\end{minipage} \\
\end{longtable}

\begin{cajaaviso}{Por qué ocupa el puesto 1 de prioridad}

Es el único hallazgo que, mal manejado, puede \textbf{destruir el activo principal de la empresa}:
la confianza de los distribuidores. Además, la decisión de abrir o no el B2C condiciona el diseño
del resto de las soluciones (digital, redes, asesoría), por eso debe resolverse primero.

\end{cajaaviso}

\needspace{6\baselineskip}

\hypertarget{modelos-inspiracionales-quiuxe9n-ya-lo-resolviuxf3}{%
\subsection*{3. Modelos inspiracionales: quién ya lo
resolvió}\label{modelos-inspiracionales-quiuxe9n-ya-lo-resolviuxf3}}

Seis empresas de sectores afines resolvieron este dilema con una misma idea: \textbf{el canal
directo complementa al distribuidor, no lo reemplaza}.

\begin{longtable}[]{@{}
  >{\raggedright\arraybackslash}p{(\columnwidth - 4\tabcolsep) * \real{0.1984}}
  >{\raggedright\arraybackslash}p{(\columnwidth - 4\tabcolsep) * \real{0.4008}}
  >{\raggedright\arraybackslash}p{(\columnwidth - 4\tabcolsep) * \real{0.4008}}@{}}
\toprule\noalign{}
\begin{minipage}[b]{\linewidth}\raggedright
\textbf{Empresa / país}
\end{minipage} & \begin{minipage}[b]{\linewidth}\raggedright
\textbf{Qué hace diferente}
\end{minipage} & \begin{minipage}[b]{\linewidth}\raggedright
\textbf{Qué puede adaptar SICHAR}
\end{minipage} \\
\midrule\noalign{}
\endfirsthead
\bottomrule\noalign{}
\endlastfoot
\textbf{Coarte SAS}

\textbf{Colombia}

\emph{\textbf{coartesas.com}} & \begin{minipage}[t]{\linewidth}\raggedright
\begin{itemize}
\item
  Trabaja con constructoras durante la obra y aprovecha su posventa para vender pequeños contratos a
  los propietarios (remodelaciones, modificaciones).
\item
  Los grandes contratos sostienen el negocio; la posventa suma ventas \textbf{sin quitarle clientes
  a la constructora}.
\end{itemize}
\end{minipage} & Acompañar a los distribuidores en la \textbf{posventa de sus proyectos} y generar
ventas complementarias con el cliente final, \textbf{en alianza con ellos y no en competencia}. \\
\textbf{LISTO Distribución}

\textbf{Colombia}

\emph{\textbf{listo.co}} & \begin{minipage}[t]{\linewidth}\raggedright
\begin{itemize}
\item
  Portafolio amplio de porcelanatos, grifería, sanitarios y accesorios para que el distribuidor
  ofrezca varias categorías.
\item
  Plataforma en línea para clientes corporativos, inventario de entrega inmediata y asesoría
  especializada.
\end{itemize}
\end{minipage} & Crear un \textbf{canal digital que promocione y genere contactos}, pero que
\textbf{dirija la compra del consumidor final hacia el distribuidor autorizado}. Más alcance sin
quitar ventas a los aliados. \\
\textbf{Roca Group}

\textbf{Internacional}

\emph{\textbf{rocagroup.com}} & \begin{minipage}[t]{\linewidth}\raggedright
\begin{itemize}
\item
  Portal B2B donde los distribuidores consultan pedidos, devoluciones, servicios y documentos en
  tiempo real.
\item
  Usa la tecnología para \textbf{fortalecer el canal tradicional}, no para reemplazarlo.
\end{itemize}
\end{minipage} & Empezar con \textbf{catálogo digital y portal o formulario para distribuidores}. Si
luego abre B2C: dirigir la compra al distribuidor más cercano o diferenciar el canal directo con
asesoría y servicios, sin los mismos productos y precios. \\
\textbf{Caterpillar}

\textbf{Internacional}

\emph{\textbf{cat.com}} & \begin{minipage}[t]{\linewidth}\raggedright
\begin{itemize}
\item
  Red global de distribuidores con repuestos rápidos y asistencia técnica continua.
\item
  Tecnología de monitoreo remoto y alto valor de reventa.
\end{itemize}
\end{minipage} & Fortalecer a la red con \textbf{soporte posventa, disponibilidad de producto y
asesoría técnica} para que el valor de la marca llegue al cliente final a través del
distribuidor. \\
\textbf{Porcelanosa}

\textbf{España}

\emph{\textbf{porcelanosa.com}} & \begin{minipage}[t]{\linewidth}\raggedright
\begin{itemize}
\item
  Combina más de 200 tiendas propias con más de 1.000 puntos de venta autorizados.
\item
  Programa «Porcelanosa Associates»: tiendas asociadas con la gama, la imagen y el respaldo de la
  marca.
\end{itemize}
\end{minipage} & \textbf{Definir qué se vende directo y qué vende la red}, y crear un programa de
\textbf{distribuidores asociados} con imagen y respaldo de la marca SICHAR. \\
\textbf{Aqualisa}

\textbf{Reino Unido}

\emph{\textbf{aqualisa.co.uk}} & \begin{minipage}[t]{\linewidth}\raggedright
\begin{itemize}
\item
  Referente en la relación entre fabricante, distribuidores e instaladores.
\item
  Se acerca al consumidor final sin perder el valor de los intermediarios.
\end{itemize}
\end{minipage} & Estrategia \textbf{omnicanal}: vender directo manteniendo a los distribuidores como
aliados mediante \textbf{comisiones, referidos, garantías, instalación o beneficios por
proyecto}. \\
\end{longtable}

\emph{Nota: los precios y cifras de terceros son referenciales, provienen de fuentes públicas de
cada empresa y deben verificarse antes de usarse en decisiones comerciales.}

\needspace{6\baselineskip}

\hypertarget{soluciuxf3n-tangible-para-sichar}{%
\subsection*{4. Solución tangible para SICHAR}\label{soluciuxf3n-tangible-para-sichar}}

La solución no es ``B2C sí o no'', sino \textbf{definir las reglas del juego antes de abrir la
puerta}. Proponemos el modelo \textbf{``Vitrina SICHAR + Distribuidor Aliado''}, que se construye en
cuatro piezas:

\needspace{12\baselineskip}

\hypertarget{pieza-1-la-poluxedtica-de-canal-un-documento-de-una-puxe1gina}{%
\subsubsection*{Pieza 1 · La política de canal (un documento de una
página)}\label{pieza-1-la-poluxedtica-de-canal-un-documento-de-una-puxe1gina}}

\begin{longtable}[]{@{}
  >{\raggedright\arraybackslash}p{(\columnwidth - 2\tabcolsep) * \real{0.2480}}
  >{\raggedright\arraybackslash}p{(\columnwidth - 2\tabcolsep) * \real{0.7520}}@{}}
\toprule\noalign{}
\begin{minipage}[b]{\linewidth}\raggedright
\textbf{Regla}
\end{minipage} & \begin{minipage}[b]{\linewidth}\raggedright
\textbf{Qué dice exactamente}
\end{minipage} \\
\midrule\noalign{}
\endfirsthead
\bottomrule\noalign{}
\endlastfoot
\textbf{Quién vende qué} & \begin{minipage}[t]{\linewidth}\raggedright
\begin{itemize}
\item
  \textbf{Línea Distribuidor:} el portafolio actual (incluido el porcelanato brillante 60x130) se
  vende solo a través de ferreterías y distribuidores.
\item
  \textbf{Línea Directa:} referencias, formatos o colecciones de proyecto exclusivos para venta
  directa (por ejemplo, kits de muestras, colecciones para arquitectos, paquetes de proyecto).
\end{itemize}
\end{minipage} \\
\textbf{Precios} & SICHAR publica un \textbf{precio sugerido de venta al público} y \textbf{no vende
directo por debajo de él}. Así el distribuidor nunca es subcotizado por su propio proveedor. \\
\textbf{Quién atiende al consumidor final} & El sitio y las redes de SICHAR \textbf{inspiran y
captan interesados}; el botón principal es \textbf{``Dónde comprar''}, que lleva al distribuidor
autorizado más cercano. El contacto se le entrega al distribuidor como \textbf{referido}. \\
\textbf{Zonas y proyectos} & Si un distribuidor ya trabaja con una constructora o proyecto, SICHAR
\textbf{no ofrece directo} a ese cliente y se registra el proyecto a nombre del distribuidor. \\
\end{longtable}

\needspace{6\baselineskip}

\hypertarget{pieza-2-el-programa-distribuidor-aliado-sichar}{%
\subsubsection*{Pieza 2 · El programa ``Distribuidor Aliado
SICHAR''}\label{pieza-2-el-programa-distribuidor-aliado-sichar}}

Un paquete de beneficios que haga que el canal \textbf{gane más} con la nueva estrategia, no menos:

\begin{itemize}
\item
  \textbf{Referidos garantizados:} los interesados captados en digital se canalizan al aliado de su
  zona.
\item
  \textbf{Muestrarios y material de punto de venta} con la imagen SICHAR (inspirado en Porcelanosa
  Associates).
\item
  \textbf{Capacitación técnica} mensual al equipo del distribuidor sobre producto y acabados.
\item
  \textbf{Protección de proyectos} registrados y condiciones preferenciales por volumen.
\item
  \textbf{Apoyo en posventa/remodelación} a clientes finales de sus proyectos (modelo Coarte).
\end{itemize}

\needspace{12\baselineskip}

\hypertarget{pieza-3-plan-de-implementaciuxf3n-por-fases}{%
\subsubsection*{Pieza 3 · Plan de implementación por
fases}\label{pieza-3-plan-de-implementaciuxf3n-por-fases}}

\begin{longtable}[]{@{}
  >{\raggedright\arraybackslash}p{(\columnwidth - 8\tabcolsep) * \real{0.0437}}
  >{\raggedright\arraybackslash}p{(\columnwidth - 8\tabcolsep) * \real{0.4405}}
  >{\raggedright\arraybackslash}p{(\columnwidth - 8\tabcolsep) * \real{0.1637}}
  >{\raggedright\arraybackslash}p{(\columnwidth - 8\tabcolsep) * \real{0.1042}}
  >{\raggedright\arraybackslash}p{(\columnwidth - 8\tabcolsep) * \real{0.2480}}@{}}
\toprule\noalign{}
\begin{minipage}[b]{\linewidth}\raggedright
\textbf{\#}
\end{minipage} & \begin{minipage}[b]{\linewidth}\raggedright
\textbf{Acción concreta}
\end{minipage} & \begin{minipage}[b]{\linewidth}\raggedright
\textbf{Responsable}
\end{minipage} & \begin{minipage}[b]{\linewidth}\raggedright
\textbf{Plazo}
\end{minipage} & \begin{minipage}[b]{\linewidth}\raggedright
\textbf{Entregable}
\end{minipage} \\
\midrule\noalign{}
\endfirsthead
\bottomrule\noalign{}
\endlastfoot
1 & Mapear el canal: lista de distribuidores y ferreterías por sede, ventas por cliente y peso de
los 20 mayores en la facturación. & Dirección comercial & Semanas 1--2 & Mapa del canal y ranking de
clientes \\
2 & Entrevistar a 8--10 distribuidores clave (mínimo 2 por sede): ¿qué temen?, ¿qué les ayudaría?,
¿qué esperan de SICHAR si se abre un canal digital? & Gerencia + coordinadores de sede & Semanas
2--4 & Informe de percepción del canal \\
3 & Redactar y validar la \textbf{política de canal} y el esquema del programa ``Distribuidor
Aliado''. & Gerencia general & Semanas 3--5 & Política de canal firmada \\
4 & Comunicar la política a \textbf{todos} los distribuidores (reunión por sede + carta),
\textbf{antes} de que exista cualquier página de venta. & Gerencia + vendedores & Semana 6 & Carta y
reunión realizadas \\
5 & Lanzar la \textbf{``Vitrina''}: catálogo digital con localizador ``Dónde comprar'' y formulario
que reparte referidos. & Líder digital & Meses 2--3 & Vitrina publicada y primeros referidos
entregados \\
6 & Evaluar con datos si conviene abrir una \textbf{Línea Directa} limitada (solo si se cumple el
semáforo de la pieza 4). & Gerencia general & Mes 6 & Decisión documentada \\
\end{longtable}

\needspace{12\baselineskip}

\hypertarget{pieza-4-el-semuxe1foro-que-decide-si-se-abre-la-venta-directa}{%
\subsubsection*{Pieza 4 · El semáforo que decide si se abre la venta
directa}\label{pieza-4-el-semuxe1foro-que-decide-si-se-abre-la-venta-directa}}

\begin{longtable}[]{@{}
  >{\raggedright\arraybackslash}p{(\columnwidth - 4\tabcolsep) * \real{0.1488}}
  >{\raggedright\arraybackslash}p{(\columnwidth - 4\tabcolsep) * \real{0.4464}}
  >{\raggedright\arraybackslash}p{(\columnwidth - 4\tabcolsep) * \real{0.4048}}@{}}
\toprule\noalign{}
\begin{minipage}[b]{\linewidth}\raggedright
\textbf{Semáforo}
\end{minipage} & \begin{minipage}[b]{\linewidth}\raggedright
\textbf{Condición}
\end{minipage} & \begin{minipage}[b]{\linewidth}\raggedright
\textbf{Decisión}
\end{minipage} \\
\midrule\noalign{}
\endfirsthead
\bottomrule\noalign{}
\endlastfoot
\textbf{VERDE} & Ningún distribuidor clave ha reducido pedidos o presentado quejas de canal, y la
Vitrina entrega referidos que se convierten en ventas para los aliados. & Se puede abrir la
\textbf{Línea Directa} limitada (colecciones exclusivas y venta a proyectos que no atiende el
canal). \\
\textbf{AMARILLO} & Hay quejas puntuales o dudas de distribuidores específicos. & Se mantiene la
Vitrina, se ajustan beneficios del programa y se vuelve a evaluar en 90 días. \\
\textbf{ROJO} & Caída sostenida de pedidos de distribuidores clave o conflicto abierto por precios.
& Se \textbf{congela} cualquier venta directa y se refuerza el programa de aliados. \\
\end{longtable}

\needspace{12\baselineskip}

\hypertarget{cuxf3mo-medir-si-estuxe1-funcionando}{%
\subsection*{5. Cómo medir si está funcionando}\label{cuxf3mo-medir-si-estuxe1-funcionando}}

\begin{longtable}[]{@{}
  >{\raggedright\arraybackslash}p{(\columnwidth - 4\tabcolsep) * \real{0.3175}}
  >{\raggedright\arraybackslash}p{(\columnwidth - 4\tabcolsep) * \real{0.3452}}
  >{\raggedright\arraybackslash}p{(\columnwidth - 4\tabcolsep) * \real{0.3373}}@{}}
\toprule\noalign{}
\begin{minipage}[b]{\linewidth}\raggedright
\textbf{Indicador}
\end{minipage} & \begin{minipage}[b]{\linewidth}\raggedright
\textbf{Cómo se mide}
\end{minipage} & \begin{minipage}[b]{\linewidth}\raggedright
\textbf{Meta sugerida (a validar con la línea base)}
\end{minipage} \\
\midrule\noalign{}
\endfirsthead
\bottomrule\noalign{}
\endlastfoot
\textbf{Retención de distribuidores clave} & Distribuidores del top 20 que siguen comprando vs.
periodo anterior & ≥ 95 \% de retención durante los primeros 12 meses \\
\textbf{Quejas por conflicto de canal} & Registro de quejas formales o informales por sede & 0
quejas críticas; resolución de las menores en \textless{} 15 días \\
\textbf{Ventas generadas por aliados} & Pedidos originados en referidos de la Vitrina & Creciente
cada trimestre desde el mes 3 \\
\textbf{\% de pedidos gestionados digitalmente} & Pedidos por formulario/portal vs. total & Línea
base en mes 1; subir progresivamente \\
\end{longtable}

\begin{cajaalerta}{Precauciones y errores que debe evitar}

\begin{itemize}
\item
  \textbf{No abrir una tienda en línea con los mismos productos y precios} antes de hablar con los
  distribuidores.
\item
  No enterarse el canal por redes sociales o por un cliente: la comunicación debe ser
  \textbf{directa y anticipada}.
\item
  No usar descuentos directos agresivos para ``probar'' el B2C: es la forma más rápida de perder la
  confianza del canal.
\end{itemize}

\end{cajaalerta}

\hallazgo{2}{\textbf{Exposición a alzas en costos de importación}}{Exposición a alzas en costos de importación}{CRÍTICA}{\textbf{Inmediato (0--30 días para el diagnóstico)}}{\textbf{Compras / Importaciones + Gerencia financiera}}

\needspace{6\baselineskip}

\hypertarget{descripciuxf3n-detallada-del-hallazgo-1}{%
\subsection*{1. Descripción detallada del hallazgo}\label{descripciuxf3n-detallada-del-hallazgo-1}}

SICHAR importa sus productos de \textbf{cinco países} (China, India, Brasil, España y Perú), con una
red de proveedores directa y sin intermediarios de importación, lo cual es una fortaleza. Sin
embargo, su estructura de costos depende de las \textbf{condiciones comerciales y los precios que
fijan esos proveedores}.

Hoy la empresa enfrenta un \textbf{incremento del 30 \% en los productos provenientes de India}.
Cada vez que sube el costo de adquisición, SICHAR debe asumir un mayor gasto para mantener la misma
oferta, y existe el riesgo adicional de que los proveedores \textbf{no mantengan las condiciones de
negociación originales}.

Si la situación no se gestiona, el efecto es directo sobre el \textbf{margen de utilidad}: la
empresa vende lo mismo, pero gana menos, o debe subir precios y arriesgarse a perder competitividad
frente a quienes tengan menores costos de adquisición.

\needspace{12\baselineskip}

\hypertarget{impacto-que-tiene-y-que-podruxeda-tener-1}{%
\subsection*{2. Impacto que tiene y que podría
tener}\label{impacto-que-tiene-y-que-podruxeda-tener-1}}

\begin{longtable}[]{@{}
  >{\raggedright\arraybackslash}p{(\columnwidth - 2\tabcolsep) * \real{0.5000}}
  >{\raggedright\arraybackslash}p{(\columnwidth - 2\tabcolsep) * \real{0.5000}}@{}}
\toprule\noalign{}
\begin{minipage}[b]{\linewidth}\raggedright
\textbf{Impacto actual (lo que ya está pasando)}
\end{minipage} & \begin{minipage}[b]{\linewidth}\raggedright
\textbf{Impacto potencial (si no se actúa)}
\end{minipage} \\
\midrule\noalign{}
\endfirsthead
\bottomrule\noalign{}
\endlastfoot
\begin{minipage}[t]{\linewidth}\raggedright
\begin{itemize}
\item
  Presión inmediata sobre los márgenes en las referencias que vienen de India.
\item
  Incertidumbre sobre precios futuros y sobre la continuidad de las condiciones pactadas con
  proveedores.
\end{itemize}
\end{minipage} & \begin{minipage}[t]{\linewidth}\raggedright
\begin{itemize}
\item
  \textbf{Reducción sostenida del margen} y menor capacidad de reinvertir en crecimiento (digital,
  asesoría, nuevos segmentos).
\item
  \textbf{Pérdida de competitividad} frente a competidores con menor costo de adquisición, como
  Cerámica Italia.
\item
  Riesgo financiero si el inventario se compra a costo alto y el mercado no acepta el ajuste de
  precio.
\end{itemize}
\end{minipage} \\
\end{longtable}

\begin{cajaaviso}{Por qué ocupa el puesto 2 de prioridad}

Es un problema \textbf{ya en marcha y con efecto directo en caja}. Además, financia todo lo demás:
sin margen sano no hay presupuesto para digitalizar, asesorar ni invertir en marca. No bloquea la
decisión del canal (hallazgo 1), pero compite con ella en urgencia.

\end{cajaaviso}

\needspace{6\baselineskip}

\hypertarget{modelos-inspiracionales-quiuxe9n-ya-lo-resolviuxf3-1}{%
\subsection*{3. Modelos inspiracionales: quién ya lo
resolvió}\label{modelos-inspiracionales-quiuxe9n-ya-lo-resolviuxf3-1}}

Cinco referentes muestran la misma lección: \textbf{no depender de un solo origen y conocer el costo
total real por referencia}.

\begin{longtable}[]{@{}
  >{\raggedright\arraybackslash}p{(\columnwidth - 4\tabcolsep) * \real{0.1984}}
  >{\raggedright\arraybackslash}p{(\columnwidth - 4\tabcolsep) * \real{0.4008}}
  >{\raggedright\arraybackslash}p{(\columnwidth - 4\tabcolsep) * \real{0.4008}}@{}}
\toprule\noalign{}
\begin{minipage}[b]{\linewidth}\raggedright
\textbf{Empresa / país}
\end{minipage} & \begin{minipage}[b]{\linewidth}\raggedright
\textbf{Qué hace diferente}
\end{minipage} & \begin{minipage}[b]{\linewidth}\raggedright
\textbf{Qué puede adaptar SICHAR}
\end{minipage} \\
\midrule\noalign{}
\endfirsthead
\bottomrule\noalign{}
\endlastfoot
\textbf{SAMISA Group}

\textbf{Internacional}

\emph{\textbf{samisa.co}} & \begin{minipage}[t]{\linewidth}\raggedright
\begin{itemize}
\item
  Conecta empresas colombianas con proveedores internacionales, principalmente en Asia.
\item
  Acompaña compra internacional y coordinación logística: optimización de costos y calidad.
\end{itemize}
\end{minipage} & \textbf{Comparar periódicamente} el costo total de importar desde China, India,
Brasil, España y Perú; negociar compras consolidadas y evaluar cada origen por costo, calidad y
tiempo. \\
\textbf{IKEA}

\textbf{Internacional}

\emph{\textbf{ikea.com}} & \begin{minipage}[t]{\linewidth}\raggedright
\begin{itemize}
\item
  Red internacional de proveedores con relaciones de largo plazo.
\item
  Presenta el abastecimiento como colaboración para lograr calidad, volumen y costo adecuados.
\end{itemize}
\end{minipage} & Calcular el \textbf{costo total puesto en Colombia por referencia}, identificar
productos que dependen de India y evaluar alternativos. Negociar \textbf{escalas de volumen,
calendarios de compra y revisión de precios}. \\
\textbf{Adidas}

\textbf{Internacional}

\emph{\textbf{adidas-group.com}} & \begin{minipage}[t]{\linewidth}\raggedright
\begin{itemize}
\item
  Diversifica su manufactura entre varios países de Asia en lugar de depender de uno solo.
\item
  Reparte así el riesgo de aranceles, tasa de cambio y costos laborales.
\end{itemize}
\end{minipage} & \textbf{Distribuir las compras de cada línea entre varios países de origen} para no
depender de uno (como India). \\
\textbf{Floor \& Decor}

\textbf{Estados Unidos}

\emph{\textbf{flooranddecor.com}} & \begin{minipage}[t]{\linewidth}\raggedright
\begin{itemize}
\item
  Compra a más de 240 proveedores en 26 países.
\item
  Redujo su dependencia de China de cerca del 50 \% (2018) a cerca del 18 \% (2024).
\item
  Mitiga aranceles negociando costos, ajustando precios y cambiando el país de origen.
\end{itemize}
\end{minipage} & \textbf{Diversificar países de origen} y mantener \textbf{más de un proveedor por
línea} de producto. \\
\textbf{Gano Excel}

\textbf{Internacional (América)}

\emph{\textbf{fuente no registrada}} & \begin{minipage}[t]{\linewidth}\raggedright
\begin{itemize}
\item
  Vende productos más costosos que el promedio mediante network marketing.
\item
  Sus distribuidores identifican el perfil de cliente y aprovechan la confianza de sus contactos.
\end{itemize}
\end{minipage} & Aprovechar la relación con los \textbf{distribuidores actuales} y acompañarlos en
su proceso comercial para sostener el \textbf{valor percibido} de los productos aunque suban los
costos. \\
\end{longtable}

\emph{Nota: los precios y cifras de terceros son referenciales, provienen de fuentes públicas de
cada empresa y deben verificarse antes de usarse en decisiones comerciales.}

\needspace{6\baselineskip}

\hypertarget{soluciuxf3n-tangible-para-sichar-1}{%
\subsection*{4. Solución tangible para SICHAR}\label{soluciuxf3n-tangible-para-sichar-1}}

La solución combina \textbf{tres palancas} que, juntas, protegen el margen: conocer el costo real,
reducir la dependencia de India y ajustar precios con método.

\needspace{6\baselineskip}

\hypertarget{ejemplo-ilustrativo-quuxe9-hace-un-30-en-el-origen-sobre-el-margen}{%
\subsubsection*{Ejemplo ilustrativo: qué hace un +30 \% en el origen sobre el
margen}\label{ejemplo-ilustrativo-quuxe9-hace-un-30-en-el-origen-sobre-el-margen}}

\emph{Cifras hipotéticas para entender la mecánica; SICHAR debe reemplazarlas por sus datos reales
en el paso 1.}

\begin{longtable}[]{@{}
  >{\raggedright\arraybackslash}p{(\columnwidth - 6\tabcolsep) * \real{0.3671}}
  >{\raggedright\arraybackslash}p{(\columnwidth - 6\tabcolsep) * \real{0.1488}}
  >{\raggedright\arraybackslash}p{(\columnwidth - 6\tabcolsep) * \real{0.2659}}
  >{\raggedright\arraybackslash}p{(\columnwidth - 6\tabcolsep) * \real{0.2183}}@{}}
\toprule\noalign{}
\begin{minipage}[b]{\linewidth}\raggedright
\textbf{Concepto}
\end{minipage} & \begin{minipage}[b]{\linewidth}\raggedright
\textbf{Antes}
\end{minipage} & \begin{minipage}[b]{\linewidth}\raggedright
\textbf{Después del +30 \% en India}
\end{minipage} & \begin{minipage}[b]{\linewidth}\raggedright
\textbf{Efecto}
\end{minipage} \\
\midrule\noalign{}
\endfirsthead
\bottomrule\noalign{}
\endlastfoot
Costo de la mercancía en origen & \$70 & \$91 (+30 \%) & +\$21 \\
Flete, seguro, aranceles y nacionalización & \$30 & \$30 & --- \\
\textbf{Costo total puesto en bodega} & \textbf{\$100} & \textbf{\$121} & \textbf{+21 \%} \\
Precio de venta al distribuidor & \$150 & \$150 & --- \\
\textbf{Margen bruto} & \textbf{33 \%} & \textbf{19 \%} & \textbf{−14 puntos} \\
Precio necesario para volver a 33 \% de margen & --- & \$181 (+21 \%) & Ajuste de precio que habría
que evaluar \\
\end{longtable}

\begin{cajaaviso}{Lectura clave}

Un alza del 30 \% en origen \textbf{no} es un alza del 30 \% en el costo total (los fletes y
aranceles no suben igual), pero \textbf{sí} puede borrar cerca de la mitad del margen. Por eso el
primer paso es medir con datos reales, referencia por referencia.

\end{cajaaviso}

\needspace{12\baselineskip}

\hypertarget{plan-de-acciuxf3n}{%
\subsubsection*{Plan de acción}\label{plan-de-acciuxf3n}}

\begin{longtable}[]{@{}
  >{\raggedright\arraybackslash}p{(\columnwidth - 8\tabcolsep) * \real{0.0437}}
  >{\raggedright\arraybackslash}p{(\columnwidth - 8\tabcolsep) * \real{0.4405}}
  >{\raggedright\arraybackslash}p{(\columnwidth - 8\tabcolsep) * \real{0.1637}}
  >{\raggedright\arraybackslash}p{(\columnwidth - 8\tabcolsep) * \real{0.1042}}
  >{\raggedright\arraybackslash}p{(\columnwidth - 8\tabcolsep) * \real{0.2480}}@{}}
\toprule\noalign{}
\begin{minipage}[b]{\linewidth}\raggedright
\textbf{\#}
\end{minipage} & \begin{minipage}[b]{\linewidth}\raggedright
\textbf{Acción concreta}
\end{minipage} & \begin{minipage}[b]{\linewidth}\raggedright
\textbf{Responsable}
\end{minipage} & \begin{minipage}[b]{\linewidth}\raggedright
\textbf{Plazo}
\end{minipage} & \begin{minipage}[b]{\linewidth}\raggedright
\textbf{Entregable}
\end{minipage} \\
\midrule\noalign{}
\endfirsthead
\bottomrule\noalign{}
\endlastfoot
1 & Construir la \textbf{hoja de costo puesto en bodega} por referencia: valor en origen, flete,
seguro, aranceles, gastos de nacionalización, almacenamiento y tasa de cambio. Empezar por las 20
referencias de mayor venta. & Compras + Finanzas & Semanas 1--3 & Hoja de costos con margen por
referencia \\
2 & Medir la \textbf{concentración por país}: qué \% de las compras y del margen depende de India
por línea de producto. Marcar en rojo las referencias donde India sea el único origen. & Compras &
Semana 3 & Mapa de dependencia por origen \\
3 & Pedir \textbf{cotizaciones y muestras} a proveedores alternativos en los países donde SICHAR ya
tiene contacto (China, Brasil, España, Perú) para las referencias en rojo; mínimo 2 opciones por
referencia. & Compras & Semanas 3--8 & ≥ 2 proveedores validados para las referencias críticas \\
4 & Renegociar con los proveedores de India: \textbf{compras consolidadas por volumen, calendario de
pedidos, pago por fases y cláusula de revisión de precios} con tope de variación. & Gerencia +
Compras & Semanas 4--10 & Acuerdo escrito o plan B documentado \\
5 & Definir la \textbf{política de precios por línea}: qué parte del alza se traslada, qué parte se
absorbe y qué referencias se reposicionan hacia productos de mayor margen (por ejemplo, el
porcelanato brillante 60x130, producto insignia). & Dirección comercial + Finanzas & Semanas 6--10 &
Lista de precios nueva por escalones \\
6 & Comunicar los ajustes a los distribuidores \textbf{con anticipación}, con argumentos claros
(costos de origen) y acompañados de beneficios (descuentos por volumen, entregas programadas). &
Dirección comercial & Semana 10--12 & Comunicado y reuniones por sede \\
7 & Consultar con la banca opciones de \textbf{cobertura cambiaria} y de compra anticipada para las
líneas más sensibles a la tasa de cambio. & Finanzas & Mes 3 & Política de cobertura aprobada \\
\end{longtable}

\needspace{12\baselineskip}

\hypertarget{cuxf3mo-medir-si-estuxe1-funcionando-1}{%
\subsection*{5. Cómo medir si está funcionando}\label{cuxf3mo-medir-si-estuxe1-funcionando-1}}

\begin{longtable}[]{@{}
  >{\raggedright\arraybackslash}p{(\columnwidth - 4\tabcolsep) * \real{0.3175}}
  >{\raggedright\arraybackslash}p{(\columnwidth - 4\tabcolsep) * \real{0.3452}}
  >{\raggedright\arraybackslash}p{(\columnwidth - 4\tabcolsep) * \real{0.3373}}@{}}
\toprule\noalign{}
\begin{minipage}[b]{\linewidth}\raggedright
\textbf{Indicador}
\end{minipage} & \begin{minipage}[b]{\linewidth}\raggedright
\textbf{Cómo se mide}
\end{minipage} & \begin{minipage}[b]{\linewidth}\raggedright
\textbf{Meta sugerida (a validar con la línea base)}
\end{minipage} \\
\midrule\noalign{}
\endfirsthead
\bottomrule\noalign{}
\endlastfoot
\textbf{Costo total de importación por referencia} & Hoja de costo puesto en bodega, actualizada
cada trimestre & 100 \% de las 20 referencias top con costo calculado en el mes 1 \\
\textbf{Margen bruto por línea} & Ventas menos costo puesto, por línea de producto & Recuperar al
menos la mitad del margen perdido en 6 meses \\
\textbf{\% de compras concentrado en India} & Compras India / compras totales por línea & Reducirlo
gradualmente; ninguna línea crítica con un solo origen \\
\textbf{Referencias con proveedor alternativo validado} & N.° de referencias críticas con ≥ 2
proveedores & 100 \% de las referencias en rojo en 12 meses \\
\end{longtable}

\begin{cajaalerta}{Precauciones y errores que debe evitar}

\begin{itemize}
\item
  Diversificar \textbf{no es solo buscar el precio más bajo}: se debe evaluar calidad, tiempos de
  entrega y costo total. Un proveedor barato con mala calidad sale más caro.
\item
  No trasladar el alza completa de golpe a los distribuidores: hágalo por escalones y con
  anticipación.
\item
  IKEA y Floor \& Decor operan a una escala mucho mayor; se toman como inspiración de método, no de
  tamaño.
\end{itemize}

\end{cajaalerta}

\hallazgo{3}{\textbf{Buyer persona no definido}}{Buyer persona no definido}{ALTA}{\textbf{Primeros 60 días (es la base de las demás soluciones)}}{\textbf{Dirección comercial + coordinadores de sede}}

\needspace{6\baselineskip}

\hypertarget{descripciuxf3n-detallada-del-hallazgo-2}{%
\subsection*{1. Descripción detallada del hallazgo}\label{descripciuxf3n-detallada-del-hallazgo-2}}

SICHAR \textbf{no tiene un registro confiable de los perfiles de sus clientes}: no se sabe con
precisión qué compran, con qué frecuencia, qué valoran, qué los frena ni cómo toman la decisión. La
relación comercial se apoya en la experiencia y la memoria de los vendedores de correría (tres por
sede).

Esto dificulta orientar el trabajo comercial y de marketing, tanto para \textbf{aprovechar mejor a
los clientes actuales} (ferreterías y distribuidores) como para \textbf{abrir segmentos nuevos} con
alto potencial: \textbf{firmas de arquitectos, constructoras y empresas de proyectos}.

Un buyer persona no es un dato demográfico: es una \textbf{radiografía de quién decide la compra,
qué problema tiene en su obra o negocio y qué lo convence}. Sin él, cualquier campaña, catálogo o
servicio se diseña a ciegas.

\needspace{12\baselineskip}

\hypertarget{impacto-que-tiene-y-que-podruxeda-tener-2}{%
\subsection*{2. Impacto que tiene y que podría
tener}\label{impacto-que-tiene-y-que-podruxeda-tener-2}}

\begin{longtable}[]{@{}
  >{\raggedright\arraybackslash}p{(\columnwidth - 2\tabcolsep) * \real{0.5000}}
  >{\raggedright\arraybackslash}p{(\columnwidth - 2\tabcolsep) * \real{0.5000}}@{}}
\toprule\noalign{}
\begin{minipage}[b]{\linewidth}\raggedright
\textbf{Impacto actual (lo que ya está pasando)}
\end{minipage} & \begin{minipage}[b]{\linewidth}\raggedright
\textbf{Impacto potencial (si no se actúa)}
\end{minipage} \\
\midrule\noalign{}
\endfirsthead
\bottomrule\noalign{}
\endlastfoot
\begin{minipage}[t]{\linewidth}\raggedright
\begin{itemize}
\item
  Mensajes comerciales y catálogo iguales para clientes con necesidades muy distintas.
\item
  Los vendedores guardan el conocimiento del cliente en su cabeza: si se van, se va la información.
\end{itemize}
\end{minipage} & \begin{minipage}[t]{\linewidth}\raggedright
\begin{itemize}
\item
  \textbf{Esfuerzo desperdiciado} en marketing y visitas mal dirigidas.
\item
  \textbf{Imposibilidad de entrar con éxito} a constructoras y arquitectos, que compran con
  criterios técnicos y de cumplimiento distintos a los de una ferretería.
\item
  Decisiones de inversión digital sin saber a quién se quiere llegar.
\end{itemize}
\end{minipage} \\
\end{longtable}

\begin{cajaaviso}{Por qué ocupa el puesto 3 de prioridad}

Es el \textbf{cimiento} de los hallazgos 4, 5 y 6: la asesoría especializada, el catálogo digital y
el contenido en redes solo funcionan si se sabe a quién se le habla. Es barato, rápido y no requiere
tecnología: solo conversaciones ordenadas con clientes.

\end{cajaaviso}

\needspace{6\baselineskip}

\hypertarget{modelos-inspiracionales-quiuxe9n-ya-lo-resolviuxf3-2}{%
\subsection*{3. Modelos inspiracionales: quién ya lo
resolvió}\label{modelos-inspiracionales-quiuxe9n-ya-lo-resolviuxf3-2}}

Seis referentes definen a sus clientes \textbf{por el problema que resuelven} y no solo por datos
demográficos.

\begin{longtable}[]{@{}
  >{\raggedright\arraybackslash}p{(\columnwidth - 4\tabcolsep) * \real{0.1984}}
  >{\raggedright\arraybackslash}p{(\columnwidth - 4\tabcolsep) * \real{0.4008}}
  >{\raggedright\arraybackslash}p{(\columnwidth - 4\tabcolsep) * \real{0.4008}}@{}}
\toprule\noalign{}
\begin{minipage}[b]{\linewidth}\raggedright
\textbf{Empresa / país}
\end{minipage} & \begin{minipage}[b]{\linewidth}\raggedright
\textbf{Qué hace diferente}
\end{minipage} & \begin{minipage}[b]{\linewidth}\raggedright
\textbf{Qué puede adaptar SICHAR}
\end{minipage} \\
\midrule\noalign{}
\endfirsthead
\bottomrule\noalign{}
\endlastfoot
\textbf{Houzz}

\textbf{Estados Unidos}

\emph{\textbf{houzz.com}} & \begin{minipage}[t]{\linewidth}\raggedright
\begin{itemize}
\item
  Organiza su oferta según las necesidades de quienes diseñan, construyen o remodelan.
\item
  Permite encontrar profesionales por especialidad y usa proyectos e ideas para que el usuario
  descubra sus gustos.
\end{itemize}
\end{minipage} & Definir perfiles distintos (\textbf{distribuidores, constructoras, arquitectos y
consumidores finales}) e investigar presupuestos, necesidades, criterios de compra y canales
preferidos \textbf{antes de diseñar campañas}. \\
\textbf{HubSpot}

\textbf{Internacional}

\emph{\textbf{blog.hubspot.com}} & \begin{minipage}[t]{\linewidth}\raggedright
\begin{itemize}
\item
  Recomienda construir buyer personas con \textbf{investigación, entrevistas y patrones reales} de
  comportamiento.
\item
  Incluye objetivos, problemas, motivaciones y forma de decidir, no solo edad o cargo.
\end{itemize}
\end{minipage} & Empezar con \textbf{tres perfiles por validar} (arquitecto/diseñador, comprador de
constructora y consumidor que remodela) y entrevistar clientes y prospectos. \\
\textbf{The Home Depot}

\textbf{Estados Unidos}

\emph{\textbf{homedepot.com}} & \begin{minipage}[t]{\linewidth}\raggedright
\begin{itemize}
\item
  Define tres segmentos con ofertas distintas: DIY (hágalo usted mismo), DIFM (hágalo por mí) y Pros
  (contratistas).
\item
  Programa Pro Xtra: descuentos, crédito y seguimiento de compras para profesionales.
\end{itemize}
\end{minipage} & Definir perfiles de \textbf{distribuidor, constructora, arquitecto y consumidor
final}, con \textbf{una oferta propia} para cada uno. \\
\textbf{Procore}

\textbf{Internacional}

\emph{\textbf{procore.com}} & \begin{minipage}[t]{\linewidth}\raggedright
\begin{itemize}
\item
  Segmenta comunicación y funcionalidades por tipo de usuario: contratistas, subcontratistas,
  propietarios, arquitectos.
\item
  Construye perfiles a partir de los \textbf{problemas de cada rol en una obra}.
\end{itemize}
\end{minipage} & Construir los perfiles desde \textbf{los problemas que enfrenta cada rol en una
obra} (constructora, arquitecto, contratista, propietario). \\
\textbf{Corona}

\textbf{Colombia}

\emph{\textbf{corona.co}} & \begin{minipage}[t]{\linewidth}\raggedright
\begin{itemize}
\item
  Oferta amplia para remodelación y clientes institucionales.
\item
  Tienda virtual, red de distribución y asesoría para distintas necesidades de compra.
\end{itemize}
\end{minipage} & Distinguir entre arquitectos, diseñadores, constructoras, distribuidores,
propietarios y compradores institucionales: \textbf{cada uno requiere productos, condiciones y
mensajes distintos}. \\
\textbf{Easy Line}

\textbf{Internacional}

\emph{\textbf{easyline.com.co}} & \begin{minipage}[t]{\linewidth}\raggedright
\begin{itemize}
\item
  Cliente muy bien definido: atiende solo a restaurantes.
\item
  Modelo enfocado en tecnología para gestionar las filas en establecimientos físicos.
\end{itemize}
\end{minipage} & Invertir en \textbf{soluciones tecnológicas que ayuden a los clientes actuales}
(ferreterías y distribuidores) a sistematizar sus procesos comerciales: fideliza y puede generar
nuevos ingresos. \\
\end{longtable}

\emph{Nota: los precios y cifras de terceros son referenciales, provienen de fuentes públicas de
cada empresa y deben verificarse antes de usarse en decisiones comerciales.}

\needspace{6\baselineskip}

\hypertarget{soluciuxf3n-tangible-para-sichar-2}{%
\subsection*{4. Solución tangible para SICHAR}\label{soluciuxf3n-tangible-para-sichar-2}}

Proponemos un \textbf{proyecto de 6 semanas} con tres entregables: entrevistas, fichas de buyer
persona y una base de clientes viva.

\needspace{12\baselineskip}

\hypertarget{perfiles-iniciales-a-validar-hipuxf3tesis-de-partida}{%
\subsubsection*{Perfiles iniciales a validar (hipótesis de
partida)}\label{perfiles-iniciales-a-validar-hipuxf3tesis-de-partida}}

\begin{longtable}[]{@{}
  >{\raggedright\arraybackslash}p{(\columnwidth - 6\tabcolsep) * \real{0.1687}}
  >{\raggedright\arraybackslash}p{(\columnwidth - 6\tabcolsep) * \real{0.2778}}
  >{\raggedright\arraybackslash}p{(\columnwidth - 6\tabcolsep) * \real{0.2778}}
  >{\raggedright\arraybackslash}p{(\columnwidth - 6\tabcolsep) * \real{0.2758}}@{}}
\toprule\noalign{}
\begin{minipage}[b]{\linewidth}\raggedright
\textbf{Campo}
\end{minipage} & \begin{minipage}[b]{\linewidth}\raggedright
\textbf{A. Distribuidor / ferretería}
\end{minipage} & \begin{minipage}[b]{\linewidth}\raggedright
\textbf{B. Constructora}
\end{minipage} & \begin{minipage}[b]{\linewidth}\raggedright
\textbf{C. Arquitecto / diseñador}
\end{minipage} \\
\midrule\noalign{}
\endfirsthead
\bottomrule\noalign{}
\endlastfoot
\textbf{Quién decide} & Propietario o gerente de compras & Jefe de compras + residente de obra & El
profesional que especifica; el cliente aprueba \\
\textbf{Qué necesita} & Rotación, buen margen, entrega rápida, respaldo del proveedor & Cumplimiento
de fechas, volumen disponible, precio por proyecto & Variedad de diseños, muestras físicas, fichas
técnicas, tendencias \\
\textbf{Qué lo frena} & Que el proveedor le compita; quiebres de inventario & Faltantes y retrasos;
lotes con diferencias de tono & No poder ver ni tocar el material; falta de asesoría \\
\textbf{Dónde se informa} & Visita del vendedor, voz a voz & Referencias de otras obras,
cotizaciones & Instagram, Pinterest, showrooms, otros colegas \\
\textbf{Qué le ofreceríamos} & Programa Distribuidor Aliado & Condiciones de proyecto +
disponibilidad garantizada & Asesoría de diseño + kit de muestras \\
\end{longtable}

\emph{Son hipótesis construidas con el diagnóstico; las entrevistas deben confirmarlas, corregirlas
o descartarlas. Un cuarto perfil ---el consumidor que remodela su vivienda--- se aborda como
``público de la Vitrina digital''.}

\needspace{12\baselineskip}

\hypertarget{plan-de-acciuxf3n-1}{%
\subsubsection*{Plan de acción}\label{plan-de-acciuxf3n-1}}

\begin{longtable}[]{@{}
  >{\raggedright\arraybackslash}p{(\columnwidth - 8\tabcolsep) * \real{0.0437}}
  >{\raggedright\arraybackslash}p{(\columnwidth - 8\tabcolsep) * \real{0.4405}}
  >{\raggedright\arraybackslash}p{(\columnwidth - 8\tabcolsep) * \real{0.1637}}
  >{\raggedright\arraybackslash}p{(\columnwidth - 8\tabcolsep) * \real{0.1042}}
  >{\raggedright\arraybackslash}p{(\columnwidth - 8\tabcolsep) * \real{0.2480}}@{}}
\toprule\noalign{}
\begin{minipage}[b]{\linewidth}\raggedright
\textbf{\#}
\end{minipage} & \begin{minipage}[b]{\linewidth}\raggedright
\textbf{Acción concreta}
\end{minipage} & \begin{minipage}[b]{\linewidth}\raggedright
\textbf{Responsable}
\end{minipage} & \begin{minipage}[b]{\linewidth}\raggedright
\textbf{Plazo}
\end{minipage} & \begin{minipage}[b]{\linewidth}\raggedright
\textbf{Entregable}
\end{minipage} \\
\midrule\noalign{}
\endfirsthead
\bottomrule\noalign{}
\endlastfoot
1 & Elegir \textbf{15--20 clientes y prospectos} (5--6 por perfil) y agendar entrevistas de 20--30
minutos, presenciales o por videollamada. & Coordinadores de sede & Semanas 1--2 & Lista de
entrevistados \\
2 & Aplicar una \textbf{guía única de preguntas} (ver recuadro) y grabar o transcribir las
respuestas en una plantilla compartida. & Vendedores + dirección comercial & Semanas 2--4 & 15--20
entrevistas documentadas \\
3 & Crear un \textbf{formulario único de visita} (en Excel compartido o formulario en línea): tipo
de cliente, productos que compra, frecuencia, tamaño del pedido, competencia que usa, motivo de
pérdida de una venta. & Dirección comercial & Semanas 2--3 & Formulario en uso por los 15
vendedores \\
4 & Analizar patrones y redactar \textbf{3 a 4 fichas de buyer persona} validadas, con mensajes y
ofertas recomendadas por perfil. & Dirección comercial & Semanas 5--6 & Fichas de buyer persona \\
5 & Capacitar a los vendedores en las fichas y en el uso del formulario; revisar los datos
\textbf{cada mes} en comité comercial. & Dirección comercial & Semana 6 en adelante & Reunión
mensual de análisis \\
\end{longtable}

\begin{cajainfo}{Guía de preguntas para las entrevistas (ejemplo)}

\begin{itemize}
\item
  ¿Cómo decide hoy qué acabados comprar y a quién? ¿Quién más participa en la decisión?
\item
  ¿Qué fue lo último que dejó de comprar a un proveedor y por qué?
\item
  ¿Qué información necesita antes de decidir (fichas técnicas, muestras, tiempos de entrega)?
\item
  ¿Dónde busca ideas o proveedores: redes, ferias, recomendaciones, catálogos?
\item
  ¿Qué haría que cambiara de proveedor? ¿Qué haría que se quedara aunque suba el precio?
\end{itemize}

\end{cajainfo}

\needspace{12\baselineskip}

\hypertarget{cuxf3mo-medir-si-estuxe1-funcionando-2}{%
\subsection*{5. Cómo medir si está funcionando}\label{cuxf3mo-medir-si-estuxe1-funcionando-2}}

\begin{longtable}[]{@{}
  >{\raggedright\arraybackslash}p{(\columnwidth - 4\tabcolsep) * \real{0.3175}}
  >{\raggedright\arraybackslash}p{(\columnwidth - 4\tabcolsep) * \real{0.3452}}
  >{\raggedright\arraybackslash}p{(\columnwidth - 4\tabcolsep) * \real{0.3373}}@{}}
\toprule\noalign{}
\begin{minipage}[b]{\linewidth}\raggedright
\textbf{Indicador}
\end{minipage} & \begin{minipage}[b]{\linewidth}\raggedright
\textbf{Cómo se mide}
\end{minipage} & \begin{minipage}[b]{\linewidth}\raggedright
\textbf{Meta sugerida (a validar con la línea base)}
\end{minipage} \\
\midrule\noalign{}
\endfirsthead
\bottomrule\noalign{}
\endlastfoot
\textbf{Entrevistas realizadas por segmento} & Conteo en la plantilla compartida & ≥ 5 por perfil en
las primeras 4 semanas \\
\textbf{Conversión de prospectos por perfil} & Cotizaciones convertidas en pedidos / cotizaciones
enviadas & Medición desde el mes 2; mejorar trimestre a trimestre \\
\textbf{Motivos de pérdida de ventas} & Registro en el formulario de visita & 100 \% de ventas
perdidas con motivo registrado \\
\textbf{Referencias más consultadas por segmento} & Consultas/cotizaciones por referencia y tipo de
cliente & Top 10 por segmento identificado en el mes 3 \\
\end{longtable}

\begin{cajaalerta}{Precauciones y errores que debe evitar}

\begin{itemize}
\item
  No dar por confirmados los perfiles iniciales: son \textbf{hipótesis} hasta que las entrevistas
  las validen.
\item
  No delegar todo en una encuesta: en este sector la información valiosa sale de
  \textbf{conversaciones con clientes}.
\end{itemize}

\end{cajaalerta}

\hallazgo{4}{\textbf{Baja madurez tecnológica y digital}}{Baja madurez tecnológica y digital}{ALTA}{\textbf{Iniciar en el mes 1; ejecutar por etapas durante 12 meses}}{\textbf{Gerencia general + líder digital (a designar)}}

\needspace{6\baselineskip}

\hypertarget{descripciuxf3n-detallada-del-hallazgo-3}{%
\subsection*{1. Descripción detallada del hallazgo}\label{descripciuxf3n-detallada-del-hallazgo-3}}

SICHAR lleva \textbf{más de 20 años operando sin una estrategia digital consolidada}. El crecimiento
se ha sostenido casi exclusivamente en el voz a voz y en una fuerza comercial de correrías. Esto
funcionó, pero hoy es un techo: la información de clientes, inventarios y cotizaciones no se
gestiona con herramientas que permitan escalar.

La falta de herramientas digitales limita la \textbf{eficiencia operativa}, dificulta
\textbf{identificar oportunidades de crecimiento} y obstaculiza la entrada a segmentos como
arquitectos y constructoras, que esperan catálogos digitales, cotizaciones rápidas y trazabilidad de
pedidos.

La buena noticia: \textbf{no se necesita una plataforma compleja ni una gran inversión inicial}. Los
referentes del sector digitalizaron por etapas, empezando por lo interno y midiendo antes de
invertir más.

\needspace{12\baselineskip}

\hypertarget{impacto-que-tiene-y-que-podruxeda-tener-3}{%
\subsection*{2. Impacto que tiene y que podría
tener}\label{impacto-que-tiene-y-que-podruxeda-tener-3}}

\begin{longtable}[]{@{}
  >{\raggedright\arraybackslash}p{(\columnwidth - 2\tabcolsep) * \real{0.5000}}
  >{\raggedright\arraybackslash}p{(\columnwidth - 2\tabcolsep) * \real{0.5000}}@{}}
\toprule\noalign{}
\begin{minipage}[b]{\linewidth}\raggedright
\textbf{Impacto actual (lo que ya está pasando)}
\end{minipage} & \begin{minipage}[b]{\linewidth}\raggedright
\textbf{Impacto potencial (si no se actúa)}
\end{minipage} \\
\midrule\noalign{}
\endfirsthead
\bottomrule\noalign{}
\endlastfoot
\begin{minipage}[t]{\linewidth}\raggedright
\begin{itemize}
\item
  Cotizaciones, pedidos e inventarios que dependen de procesos manuales o de la memoria de personas.
\item
  Imposibilidad de ver datos consolidados de las cinco sedes para tomar decisiones.
\end{itemize}
\end{minipage} & \begin{minipage}[t]{\linewidth}\raggedright
\begin{itemize}
\item
  \textbf{Pérdida de eficiencia} y errores de pedido a medida que crece la operación.
\item
  \textbf{Dificultad para expandirse} a nuevos segmentos que exigen respuesta rápida y catálogo
  digital.
\item
  Pérdida de terreno frente a competidores que usan la tecnología para mejorar procesos y relación
  con clientes.
\end{itemize}
\end{minipage} \\
\end{longtable}

\begin{cajaaviso}{Por qué ocupa el puesto 4 de prioridad}

Es la \textbf{infraestructura} que hace posibles los hallazgos 1, 5 y 6: la Vitrina, el localizador
``Dónde comprar'', el catálogo y las redes dependen de tener información ordenada. Su impacto se ve
en el mediano plazo, por eso va después de los riesgos que afectan caja y canal.

\end{cajaaviso}

\needspace{6\baselineskip}

\hypertarget{modelos-inspiracionales-quiuxe9n-ya-lo-resolviuxf3-3}{%
\subsection*{3. Modelos inspiracionales: quién ya lo
resolvió}\label{modelos-inspiracionales-quiuxe9n-ya-lo-resolviuxf3-3}}

Cinco referentes ---desde distribuidores industriales hasta una startup agrícola--- muestran que se
puede \textbf{empezar simple y crecer por etapas}.

\begin{longtable}[]{@{}
  >{\raggedright\arraybackslash}p{(\columnwidth - 4\tabcolsep) * \real{0.1984}}
  >{\raggedright\arraybackslash}p{(\columnwidth - 4\tabcolsep) * \real{0.4008}}
  >{\raggedright\arraybackslash}p{(\columnwidth - 4\tabcolsep) * \real{0.4008}}@{}}
\toprule\noalign{}
\begin{minipage}[b]{\linewidth}\raggedright
\textbf{Empresa / país}
\end{minipage} & \begin{minipage}[b]{\linewidth}\raggedright
\textbf{Qué hace diferente}
\end{minipage} & \begin{minipage}[b]{\linewidth}\raggedright
\textbf{Qué puede adaptar SICHAR}
\end{minipage} \\
\midrule\noalign{}
\endfirsthead
\bottomrule\noalign{}
\endlastfoot
\textbf{Grainger}

\textbf{Internacional}

\emph{\textbf{grainger.com}} & \begin{minipage}[t]{\linewidth}\raggedright
\begin{itemize}
\item
  Herramientas digitales para que clientes empresariales busquen productos, consulten información y
  gestionen compras.
\item
  Un distribuidor tradicional que facilita la compra con catálogo y servicios digitales.
\end{itemize}
\end{minipage} & \textbf{Catálogo digital ordenado por líneas} con fotos, medidas, fichas técnicas y
disponibilidad. Luego, solicitudes de cotización e historial de pedidos. Empezar con herramientas
sencillas y medir uso. \\
\textbf{Grupo Decor}

\textbf{Colombia}

\emph{\textbf{grupodecor.com}} & \begin{minipage}[t]{\linewidth}\raggedright
\begin{itemize}
\item
  Portafolio de acabados para distribuidores, constructoras y proyectos.
\item
  Diferencia la atención a distribuidores de la atención especializada a constructoras y grandes
  proyectos.
\end{itemize}
\end{minipage} & Estrategia digital con \textbf{catálogo actualizado, fotos de proyectos reales y
contenido sobre combinaciones}, que facilite contactar a SICHAR o a sus distribuidores
autorizados. \\
\textbf{Sodimac Colombia}

\textbf{Colombia}

\emph{\textbf{homecenter.com.co}} & \begin{minipage}[t]{\linewidth}\raggedright
\begin{itemize}
\item
  Opera con Corona la Aceleradora Sodimac Corona, programa de innovación abierta con startups.
\end{itemize}
\end{minipage} & Hacer \textbf{pilotos digitales con aliados tecnológicos} (startups o proveedores)
antes de grandes inversiones. \\
\textbf{Cemex}

\textbf{México / Internacional}

\emph{\textbf{cemex.com}} & \begin{minipage}[t]{\linewidth}\raggedright
\begin{itemize}
\item
  Desarrolló «CEMEX Go», plataforma B2B donde distribuidores y constructoras hacen pedidos.
\item
  Digitalizó \textbf{por etapas}: primero procesos internos, luego autoservicio y después
  comunicación.
\end{itemize}
\end{minipage} & Digitalizar por etapas: \textbf{(1) inventario y cotizaciones, (2) autoservicio
para distribuidores y constructoras, (3) comunicación y redes}. \\
\textbf{ComproAgro}

\textbf{Colombia}

\emph{\textbf{comproagro.com}} & \begin{minipage}[t]{\linewidth}\raggedright
\begin{itemize}
\item
  Plataforma que conecta directamente a campesinos con compradores, eliminando intermediarios.
\item
  Permite fijar el precio antes de transportar los productos.
\end{itemize}
\end{minipage} & Usar \textbf{bases de datos y herramientas de análisis de mercado} para anticipar y
fijar con tiempo los precios de los insumos que se importan. \\
\end{longtable}

\emph{Nota: los precios y cifras de terceros son referenciales, provienen de fuentes públicas de
cada empresa y deben verificarse antes de usarse en decisiones comerciales.}

\needspace{6\baselineskip}

\hypertarget{soluciuxf3n-tangible-para-sichar-3}{%
\subsection*{4. Solución tangible para SICHAR}\label{soluciuxf3n-tangible-para-sichar-3}}

Proponemos la \textbf{``Ruta Digital SICHAR''} en tres etapas, donde cada una solo se activa si la
anterior se está usando de verdad.

\begin{longtable}[]{@{}
  >{\raggedright\arraybackslash}p{(\columnwidth - 6\tabcolsep) * \real{0.1786}}
  >{\raggedright\arraybackslash}p{(\columnwidth - 6\tabcolsep) * \real{0.3869}}
  >{\raggedright\arraybackslash}p{(\columnwidth - 6\tabcolsep) * \real{0.2381}}
  >{\raggedright\arraybackslash}p{(\columnwidth - 6\tabcolsep) * \real{0.1964}}@{}}
\toprule\noalign{}
\begin{minipage}[b]{\linewidth}\raggedright
\textbf{Etapa}
\end{minipage} & \begin{minipage}[b]{\linewidth}\raggedright
\textbf{Qué se hace}
\end{minipage} & \begin{minipage}[b]{\linewidth}\raggedright
\textbf{Herramientas sugeridas (bajo costo)}
\end{minipage} & \begin{minipage}[b]{\linewidth}\raggedright
\textbf{Resultado visible}
\end{minipage} \\
\midrule\noalign{}
\endfirsthead
\bottomrule\noalign{}
\endlastfoot
\textbf{Etapa 1} · Ordenar la casa

(Meses 1--3) & \begin{minipage}[t]{\linewidth}\raggedright
\begin{itemize}
\item
  Un solo archivo maestro de productos (referencia, medidas, foto, precio, existencias).
\item
  Plantilla de cotización estándar y registro de cotizaciones.
\item
  Formulario de visita para los vendedores (ligado al hallazgo 3).
\item
  Catálogo digital en PDF por línea.
\end{itemize}
\end{minipage} & Hojas de cálculo compartidas en la nube, WhatsApp Business con catálogo, CRM
gratuito & Cotizar en menos de 24--48 horas con datos consistentes en las 5 sedes \\
\textbf{Etapa 2} · Abrir la puerta

(Meses 4--6) & \begin{minipage}[t]{\linewidth}\raggedright
\begin{itemize}
\item
  Sitio web sencillo con catálogo por línea y proyectos reales.
\item
  Formulario de cotización y \textbf{localizador ``Dónde comprar''}.
\item
  Pedidos de distribuidores por formulario o portal básico.
\end{itemize}
\end{minipage} & Constructor de sitios web, formularios con respuesta automática, CRM & Primeros
contactos digitales y pedidos gestionados en línea \\
\textbf{Etapa 3} · Autoservicio

(Meses 7--12) & \begin{minipage}[t]{\linewidth}\raggedright
\begin{itemize}
\item
  Portal para distribuidores y constructoras: consulta de inventario, historial y seguimiento de
  pedidos.
\item
  Tablero de indicadores por sede.
\item
  Integración con inventarios y facturación, si el piloto lo justifica.
\end{itemize}
\end{minipage} & Piloto con aliado tecnológico (modelo Sodimac/Corona) antes de comprar nada grande
& Menos llamadas operativas; datos para decidir \\
\end{longtable}

\needspace{12\baselineskip}

\hypertarget{plan-de-acciuxf3n-de-los-primeros-90-duxedas}{%
\subsubsection*{Plan de acción de los primeros 90
días}\label{plan-de-acciuxf3n-de-los-primeros-90-duxedas}}

\begin{longtable}[]{@{}
  >{\raggedright\arraybackslash}p{(\columnwidth - 8\tabcolsep) * \real{0.0437}}
  >{\raggedright\arraybackslash}p{(\columnwidth - 8\tabcolsep) * \real{0.4405}}
  >{\raggedright\arraybackslash}p{(\columnwidth - 8\tabcolsep) * \real{0.1637}}
  >{\raggedright\arraybackslash}p{(\columnwidth - 8\tabcolsep) * \real{0.1042}}
  >{\raggedright\arraybackslash}p{(\columnwidth - 8\tabcolsep) * \real{0.2480}}@{}}
\toprule\noalign{}
\begin{minipage}[b]{\linewidth}\raggedright
\textbf{\#}
\end{minipage} & \begin{minipage}[b]{\linewidth}\raggedright
\textbf{Acción concreta}
\end{minipage} & \begin{minipage}[b]{\linewidth}\raggedright
\textbf{Responsable}
\end{minipage} & \begin{minipage}[b]{\linewidth}\raggedright
\textbf{Plazo}
\end{minipage} & \begin{minipage}[b]{\linewidth}\raggedright
\textbf{Entregable}
\end{minipage} \\
\midrule\noalign{}
\endfirsthead
\bottomrule\noalign{}
\endlastfoot
1 & Designar un \textbf{líder digital} interno (puede ser medio tiempo) con autoridad para coordinar
a las 5 sedes. & Gerencia general & Semana 1 & Líder nombrado \\
2 & Hacer el \textbf{inventario de procesos}: cómo se cotiza, se pide y se factura hoy en cada sede;
detectar los 3 cuellos de botella. & Líder digital & Semanas 1--3 & Diagnóstico de procesos \\
3 & Crear el \textbf{archivo maestro de productos} y el catálogo digital de las 4 líneas (accesorios
de baño, cerámicas y acabados, elementos decorativos, remodelación). & Líder digital + comercial &
Semanas 3--8 & Catálogo digital v1 \\
4 & Implementar una herramienta de \textbf{gestión de clientes (CRM) básica} y migrar los datos del
formulario de visita. & Líder digital & Semanas 6--10 & CRM en uso en 5 sedes \\
5 & Capacitar a los 15 vendedores en el uso de las herramientas y fijar la regla: \textbf{lo que no
está registrado, no existió}. & Dirección comercial & Semanas 8--12 & Capacitación realizada \\
6 & Revisar en el mes 3 el uso real y decidir si se activa la Etapa 2. & Gerencia general & Mes 3 &
Decisión documentada \\
\end{longtable}

\needspace{12\baselineskip}

\hypertarget{cuxf3mo-medir-si-estuxe1-funcionando-3}{%
\subsection*{5. Cómo medir si está funcionando}\label{cuxf3mo-medir-si-estuxe1-funcionando-3}}

\begin{longtable}[]{@{}
  >{\raggedright\arraybackslash}p{(\columnwidth - 4\tabcolsep) * \real{0.3175}}
  >{\raggedright\arraybackslash}p{(\columnwidth - 4\tabcolsep) * \real{0.3452}}
  >{\raggedright\arraybackslash}p{(\columnwidth - 4\tabcolsep) * \real{0.3373}}@{}}
\toprule\noalign{}
\begin{minipage}[b]{\linewidth}\raggedright
\textbf{Indicador}
\end{minipage} & \begin{minipage}[b]{\linewidth}\raggedright
\textbf{Cómo se mide}
\end{minipage} & \begin{minipage}[b]{\linewidth}\raggedright
\textbf{Meta sugerida (a validar con la línea base)}
\end{minipage} \\
\midrule\noalign{}
\endfirsthead
\bottomrule\noalign{}
\endlastfoot
\textbf{\% del catálogo actualizado} & Referencias con foto, medidas, ficha y precio / total & ≥ 90
\% en el mes 3 \\
\textbf{Tiempo de respuesta a cotizaciones} & Horas entre solicitud y envío & \textless{} 48 horas,
y \textless{} 24 horas para referencias top \\
\textbf{Oportunidades registradas digitalmente} & Registros en CRM / visitas realizadas & ≥ 80 \% de
visitas registradas en el mes 4 \\
\textbf{Pedidos iniciados por canales digitales} & Pedidos por formulario/portal / total & Línea
base en la Etapa 2; crecimiento trimestral \\
\textbf{Errores de pedido} & Pedidos con error / total & Reducción trimestral frente a la línea
base \\
\end{longtable}

\begin{cajaalerta}{Precauciones y errores que debe evitar}

\begin{itemize}
\item
  \textbf{No comprar una plataforma grande primero.} Primero se ordenan datos y procesos; luego se
  automatiza.
\item
  Grainger y Cemex tienen una escala muy superior: se toman como \textbf{modelo de método}, no de
  tamaño.
\item
  La tecnología sin capacitación no se usa: reserve tiempo para entrenar a los vendedores.
\end{itemize}

\end{cajaalerta}

\hallazgo{5}{\textbf{Falta de atención especializada}}{Falta de atención especializada}{MEDIA-ALTA}{\textbf{Piloto en los meses 3--6, una vez definidos los perfiles}}{\textbf{Dirección comercial + asesor de diseño}}

\needspace{6\baselineskip}

\hypertarget{descripciuxf3n-detallada-del-hallazgo-4}{%
\subsection*{1. Descripción detallada del hallazgo}\label{descripciuxf3n-detallada-del-hallazgo-4}}

SICHAR \textbf{no cuenta con un servicio de asesoría especializada} que ayude a sus clientes a
seleccionar productos y acabados según las necesidades de cada proyecto. La venta se limita a
ofrecer producto, sin acompañamiento técnico ni de diseño.

Esta ausencia pesa especialmente con \textbf{arquitectos, diseñadores y constructoras}, que no
buscan solo un porcelanato: buscan orientación técnica, recomendaciones según presupuesto y estilo,
y respaldo para especificar con tranquilidad. Es, además, una de las \textbf{oportunidades
identificadas} en el diagnóstico: ofrecer asesoría como diferenciador frente a la competencia.

Hoy el acompañamiento personalizado lo ofrecen otros competidores, lo que deja a SICHAR compitiendo
casi solo por producto y precio, justo donde el hallazgo 2 le resta margen.

\needspace{12\baselineskip}

\hypertarget{impacto-que-tiene-y-que-podruxeda-tener-4}{%
\subsection*{2. Impacto que tiene y que podría
tener}\label{impacto-que-tiene-y-que-podruxeda-tener-4}}

\begin{longtable}[]{@{}
  >{\raggedright\arraybackslash}p{(\columnwidth - 2\tabcolsep) * \real{0.5000}}
  >{\raggedright\arraybackslash}p{(\columnwidth - 2\tabcolsep) * \real{0.5000}}@{}}
\toprule\noalign{}
\begin{minipage}[b]{\linewidth}\raggedright
\textbf{Impacto actual (lo que ya está pasando)}
\end{minipage} & \begin{minipage}[b]{\linewidth}\raggedright
\textbf{Impacto potencial (si no se actúa)}
\end{minipage} \\
\midrule\noalign{}
\endfirsthead
\bottomrule\noalign{}
\endlastfoot
\begin{minipage}[t]{\linewidth}\raggedright
\begin{itemize}
\item
  Los profesionales del diseño y la construcción no encuentran en SICHAR un socio técnico.
\item
  La venta depende solo del producto y del precio, sin un elemento diferenciador.
\end{itemize}
\end{minipage} & \begin{minipage}[t]{\linewidth}\raggedright
\begin{itemize}
\item
  \textbf{Pérdida de oportunidades comerciales} frente a competidores que sí acompañan.
\item
  Dificultad para \textbf{fidelizar} y construir relaciones duraderas con arquitectos y
  constructoras.
\item
  Menor capacidad de defender el precio ante alzas de costos.
\end{itemize}
\end{minipage} \\
\end{longtable}

\begin{cajaaviso}{Por qué ocupa el puesto 5 de prioridad}

Es el \textbf{diferenciador} que permite entrar a nuevos segmentos con un argumento distinto al
precio. Va en quinto lugar porque necesita primero conocer a esos clientes (hallazgo 3) y tener
catálogo y herramientas básicas (hallazgo 4).

\end{cajaaviso}

\needspace{6\baselineskip}

\hypertarget{modelos-inspiracionales-quiuxe9n-ya-lo-resolviuxf3-4}{%
\subsection*{3. Modelos inspiracionales: quién ya lo
resolvió}\label{modelos-inspiracionales-quiuxe9n-ya-lo-resolviuxf3-4}}

Cinco referentes compiten \textbf{con el acompañamiento y no solo con el producto}.

\begin{longtable}[]{@{}
  >{\raggedright\arraybackslash}p{(\columnwidth - 4\tabcolsep) * \real{0.2203}}
  >{\raggedright\arraybackslash}p{(\columnwidth - 4\tabcolsep) * \real{0.3900}}
  >{\raggedright\arraybackslash}p{(\columnwidth - 4\tabcolsep) * \real{0.3897}}@{}}
\toprule\noalign{}
\begin{minipage}[b]{\linewidth}\raggedright
\textbf{Empresa / país}
\end{minipage} & \begin{minipage}[b]{\linewidth}\raggedright
\textbf{Qué hace diferente}
\end{minipage} & \begin{minipage}[b]{\linewidth}\raggedright
\textbf{Qué puede adaptar SICHAR}
\end{minipage} \\
\midrule\noalign{}
\endfirsthead
\bottomrule\noalign{}
\endlastfoot
\textbf{Corona}

\textbf{Colombia}

\emph{\textbf{corona.co}} & \begin{minipage}[t]{\linewidth}\raggedright
\begin{itemize}
\item
  Centros especializados donde los clientes conocen y comparan acabados.
\item
  Asesoría personalizada y herramientas de diseño para proyectos de remodelación.
\end{itemize}
\end{minipage} & Servicio de asesoría para \textbf{arquitectos, diseñadores y constructoras} con
recomendaciones según presupuesto, estilo y necesidades del proyecto. \\
\textbf{The Home Depot}

\textbf{Estados Unidos}

\emph{\textbf{corporate.homedepot.com}} & \begin{minipage}[t]{\linewidth}\raggedright
\begin{itemize}
\item
  Consultas con especialistas, presenciales y virtuales; la consulta inicial es gratuita.
\item
  Proceso: revisión del espacio, selección de productos, diseño, presupuesto e instalación.
\end{itemize}
\end{minipage} & \textbf{Consultas de diseño} (presenciales o virtuales) con revisión del espacio,
selección de productos y presupuesto, con conexión a \textbf{instaladores aliados}. \\
\textbf{The Tile Shop}

\textbf{Estados Unidos}

\emph{\textbf{tileshop.com}} & \begin{minipage}[t]{\linewidth}\raggedright
\begin{itemize}
\item
  Destaca sus capacidades de diseño, fabricación y distribución para propietarios y profesionales.
\end{itemize}
\end{minipage} & Ofrecer un \textbf{servicio de diseño incluido con la compra}. \\
\textbf{Porcelanosa}

\textbf{España / Internacional}

\emph{\textbf{porcelanosa.com}} & \begin{minipage}[t]{\linewidth}\raggedright
\begin{itemize}
\item
  Showrooms con asesores que recomiendan acabados según estilo, presupuesto y tipo de proyecto.
\item
  Trabaja con arquitectos y diseñadores: muestras físicas y asesoría técnica para especificar.
\end{itemize}
\end{minipage} & Formar asesores y crear un \textbf{programa de muestras físicas y asesoría técnica}
para arquitectos y diseñadores. \\
\textbf{Grupo Bios}

\textbf{Colombia}

\emph{\textbf{fuente no registrada}} & \begin{minipage}[t]{\linewidth}\raggedright
\begin{itemize}
\item
  Líneas de negocio por tipo de cliente: canal tradicional atendido por vendedores y canal moderno
  por Key Account Managers.
\end{itemize}
\end{minipage} & \textbf{Organizar a los clientes por tamaño de pedido}: a mayor tamaño, mayor
especialización y nivel de atención del equipo comercial. \\
\end{longtable}

\emph{Nota: los precios y cifras de terceros son referenciales, provienen de fuentes públicas de
cada empresa y deben verificarse antes de usarse en decisiones comerciales.}

\needspace{6\baselineskip}

\hypertarget{soluciuxf3n-tangible-para-sichar-4}{%
\subsection*{4. Solución tangible para SICHAR}\label{soluciuxf3n-tangible-para-sichar-4}}

Proponemos crear la \textbf{``Mesa de Asesoría SICHAR''}: un servicio pequeño, medible y escalable
que empieza como piloto en la sede de \textbf{Cali} (origen de la empresa) y se replica si funciona.

\needspace{12\baselineskip}

\hypertarget{quuxe9-incluye-el-servicio}{%
\subsubsection*{Qué incluye el servicio}\label{quuxe9-incluye-el-servicio}}

\begin{longtable}[]{@{}
  >{\raggedright\arraybackslash}p{(\columnwidth - 4\tabcolsep) * \real{0.2679}}
  >{\raggedright\arraybackslash}p{(\columnwidth - 4\tabcolsep) * \real{0.4960}}
  >{\raggedright\arraybackslash}p{(\columnwidth - 4\tabcolsep) * \real{0.2361}}@{}}
\toprule\noalign{}
\begin{minipage}[b]{\linewidth}\raggedright
\textbf{Servicio}
\end{minipage} & \begin{minipage}[b]{\linewidth}\raggedright
\textbf{En qué consiste}
\end{minipage} & \begin{minipage}[b]{\linewidth}\raggedright
\textbf{Para quién}
\end{minipage} \\
\midrule\noalign{}
\endfirsthead
\bottomrule\noalign{}
\endlastfoot
\textbf{1. Kit de muestras por proyecto} & Caja con muestras físicas de las referencias recomendadas
para cada espacio (baño, cocina, fachada, zona social). & Arquitectos y diseñadores \\
\textbf{2. Cotización de proyecto en 48 horas} & Lista de referencias, cantidades por m², costo
estimado y tiempos de entrega, a partir de planos o medidas. & Constructoras y arquitectos \\
\textbf{3. Asesoría de 30 minutos} & Videollamada o visita con un asesor para elegir acabados según
estilo, presupuesto y tipo de proyecto. & Todos los perfiles \\
\textbf{4. Tablero de inspiración (moodboard)} & Documento visual con la combinación de pisos,
revestimientos, grifería y decorativos propuesta. & Diseñadores y propietarios \\
\textbf{5. Programa ``Arquitecto Aliado''} & Muestras gratuitas, condiciones preferenciales de
proyecto, ficha técnica a la mano y acceso prioritario a lanzamientos. & Firmas de arquitectos y
diseño \\
\end{longtable}

\begin{cajainfo}{Cómo se conecta con el hallazgo 1}

La asesoría \textbf{no vende por fuera del canal}: cuando el cliente es consumidor final o un
proyecto pequeño, el asesor entrega la cotización y \textbf{remite la compra al distribuidor aliado
de la zona}. Así el servicio refuerza la red en lugar de competir con ella.

\end{cajainfo}

\needspace{12\baselineskip}

\hypertarget{plan-de-acciuxf3n-2}{%
\subsubsection*{Plan de acción}\label{plan-de-acciuxf3n-2}}

\begin{longtable}[]{@{}
  >{\raggedright\arraybackslash}p{(\columnwidth - 8\tabcolsep) * \real{0.0437}}
  >{\raggedright\arraybackslash}p{(\columnwidth - 8\tabcolsep) * \real{0.4405}}
  >{\raggedright\arraybackslash}p{(\columnwidth - 8\tabcolsep) * \real{0.1637}}
  >{\raggedright\arraybackslash}p{(\columnwidth - 8\tabcolsep) * \real{0.1042}}
  >{\raggedright\arraybackslash}p{(\columnwidth - 8\tabcolsep) * \real{0.2480}}@{}}
\toprule\noalign{}
\begin{minipage}[b]{\linewidth}\raggedright
\textbf{\#}
\end{minipage} & \begin{minipage}[b]{\linewidth}\raggedright
\textbf{Acción concreta}
\end{minipage} & \begin{minipage}[b]{\linewidth}\raggedright
\textbf{Responsable}
\end{minipage} & \begin{minipage}[b]{\linewidth}\raggedright
\textbf{Plazo}
\end{minipage} & \begin{minipage}[b]{\linewidth}\raggedright
\textbf{Entregable}
\end{minipage} \\
\midrule\noalign{}
\endfirsthead
\bottomrule\noalign{}
\endlastfoot
1 & Definir el \textbf{perfil del asesor}: conocimiento de acabados y trato con arquitectos.
Opciones de bajo riesgo: una diseñadora o arquitecta aliada pagada por proyecto, o un asesor interno
en la sede Cali. & Gerencia + comercial & Mes 3 & Asesor o aliado seleccionado \\
2 & Armar la \textbf{guía de selección de acabados}: qué producto recomendar por tipo de espacio,
estilo y presupuesto, con las 4 líneas del portafolio (accesorios de baño, cerámicas y acabados,
elementos decorativos, remodelación). & Asesor + compras & Mes 3--4 & Guía de selección v1 \\
3 & Producir \textbf{10 kits de muestras} y las plantillas de cotización de proyecto y moodboard. &
Asesor + líder digital & Mes 4 & Kits y plantillas listos \\
4 & Lanzar el piloto invitando a \textbf{10--15 arquitectos, diseñadores y constructoras} de Cali
identificados en las entrevistas del hallazgo 3. & Dirección comercial & Mes 4 & Primeros proyectos
asesorados \\
5 & Capacitar a los \textbf{15 vendedores} en producto técnico (sesión mensual de 2 horas) para que
deriven consultas al asesor. & Asesor & Mes 4 en adelante & Calendario de capacitaciones \\
6 & Evaluar el piloto a los 90 días y decidir si se \textbf{replica en las otras 4 sedes}. &
Gerencia general & Mes 7 & Informe de resultados \\
\end{longtable}

\needspace{12\baselineskip}

\hypertarget{cuxf3mo-medir-si-estuxe1-funcionando-4}{%
\subsection*{5. Cómo medir si está funcionando}\label{cuxf3mo-medir-si-estuxe1-funcionando-4}}

\begin{longtable}[]{@{}
  >{\raggedright\arraybackslash}p{(\columnwidth - 4\tabcolsep) * \real{0.3175}}
  >{\raggedright\arraybackslash}p{(\columnwidth - 4\tabcolsep) * \real{0.3452}}
  >{\raggedright\arraybackslash}p{(\columnwidth - 4\tabcolsep) * \real{0.3373}}@{}}
\toprule\noalign{}
\begin{minipage}[b]{\linewidth}\raggedright
\textbf{Indicador}
\end{minipage} & \begin{minipage}[b]{\linewidth}\raggedright
\textbf{Cómo se mide}
\end{minipage} & \begin{minipage}[b]{\linewidth}\raggedright
\textbf{Meta sugerida (a validar con la línea base)}
\end{minipage} \\
\midrule\noalign{}
\endfirsthead
\bottomrule\noalign{}
\endlastfoot
\textbf{Proyectos asesorados} & N.° de proyectos con asesoría & ≥ 10--15 en los primeros 90 días del
piloto \\
\textbf{Cotizaciones convertidas en pedidos} & Pedidos / cotizaciones de proyecto & Mejorar frente
al promedio sin asesoría \\
\textbf{Valor promedio por proyecto} & Ventas por proyecto asesorado & Superior al ticket promedio
actual \\
\textbf{Satisfacción del cliente} & Encuesta corta posterior al proyecto & ≥ 4,5 sobre 5 \\
\textbf{Tiempo invertido por asesoría} & Horas del asesor por proyecto & Controlar para que el
servicio sea rentable \\
\end{longtable}

\begin{cajaalerta}{Precauciones y errores que debe evitar}

\begin{itemize}
\item
  El modelo de The Home Depot está pensado para su escala: SICHAR debe \textbf{ajustar el alcance a
  sus recursos} y empezar pequeño.
\item
  No ofrecer diseño ``gratis'' sin condiciones: vincúlelo a compra o a membresía del programa para
  que sea sostenible.
\end{itemize}

\end{cajaalerta}

\hallazgo{6}{\textbf{Baja presencia en redes sociales}}{Baja presencia en redes sociales}{MEDIA}{\textbf{Iniciar en el mes 3 y mantener de forma continua}}{\textbf{Líder digital + apoyo creativo}}

\needspace{6\baselineskip}

\hypertarget{descripciuxf3n-detallada-del-hallazgo-5}{%
\subsection*{1. Descripción detallada del hallazgo}\label{descripciuxf3n-detallada-del-hallazgo-5}}

SICHAR tiene una \textbf{presencia limitada en redes sociales}, lo que reduce la visibilidad de su
marca y sus productos frente a clientes potenciales. Aunque la empresa tiene reconocimiento sólido
por voz a voz, ese reconocimiento \textbf{no es visible} para quien descubre proveedores hoy, y
mucho menos para los nuevos segmentos.

Competidores como \textbf{Cerámica Italia} ya aprovechan los canales digitales para exhibir
producto, comunicar su propuesta de valor y construir reconocimiento. Los arquitectos, diseñadores y
constructoras usan las redes para \textbf{descubrir proveedores e identificar tendencias} del
sector.

Para un negocio cuyo producto es visual ---porcelanatos, revestimientos, grifería, decorativos--- no
estar presente en redes es dejar de mostrar su mejor argumento de venta.

\needspace{12\baselineskip}

\hypertarget{impacto-que-tiene-y-que-podruxeda-tener-5}{%
\subsection*{2. Impacto que tiene y que podría
tener}\label{impacto-que-tiene-y-que-podruxeda-tener-5}}

\begin{longtable}[]{@{}
  >{\raggedright\arraybackslash}p{(\columnwidth - 2\tabcolsep) * \real{0.5000}}
  >{\raggedright\arraybackslash}p{(\columnwidth - 2\tabcolsep) * \real{0.5000}}@{}}
\toprule\noalign{}
\begin{minipage}[b]{\linewidth}\raggedright
\textbf{Impacto actual (lo que ya está pasando)}
\end{minipage} & \begin{minipage}[b]{\linewidth}\raggedright
\textbf{Impacto potencial (si no se actúa)}
\end{minipage} \\
\midrule\noalign{}
\endfirsthead
\bottomrule\noalign{}
\endlastfoot
\begin{minipage}[t]{\linewidth}\raggedright
\begin{itemize}
\item
  Baja visibilidad de marca y de producto insignia fuera del círculo de clientes actuales.
\item
  Sin canal para mostrar proyectos reales ni tendencias.
\end{itemize}
\end{minipage} & \begin{minipage}[t]{\linewidth}\raggedright
\begin{itemize}
\item
  \textbf{Dificultad para atraer nuevos clientes} y posicionar la marca.
\item
  \textbf{Pérdida de oportunidades comerciales} con arquitectos, diseñadores y constructoras que
  buscan proveedores en redes.
\item
  Ventaja creciente de competidores más visibles.
\end{itemize}
\end{minipage} \\
\end{longtable}

\begin{cajaaviso}{Por qué ocupa el puesto 6 de prioridad}

Es el hallazgo \textbf{menos urgente} porque depende de los anteriores: sin política de canal, buyer
persona y catálogo, publicar produce ruido y no ventas. Pero es de bajo costo y alta visibilidad,
por lo que conviene empezar en el mes 3 y sostenerlo.

\end{cajaaviso}

\needspace{6\baselineskip}

\hypertarget{modelos-inspiracionales-quiuxe9n-ya-lo-resolviuxf3-5}{%
\subsection*{3. Modelos inspiracionales: quién ya lo
resolvió}\label{modelos-inspiracionales-quiuxe9n-ya-lo-resolviuxf3-5}}

Los referentes del sector usan lo digital para \textbf{captar interesados y llevarlos a un punto de
contacto}, no solo para publicar.

\begin{longtable}[]{@{}
  >{\raggedright\arraybackslash}p{(\columnwidth - 4\tabcolsep) * \real{0.1984}}
  >{\raggedright\arraybackslash}p{(\columnwidth - 4\tabcolsep) * \real{0.4008}}
  >{\raggedright\arraybackslash}p{(\columnwidth - 4\tabcolsep) * \real{0.4008}}@{}}
\toprule\noalign{}
\begin{minipage}[b]{\linewidth}\raggedright
\textbf{Empresa / país}
\end{minipage} & \begin{minipage}[b]{\linewidth}\raggedright
\textbf{Qué hace diferente}
\end{minipage} & \begin{minipage}[b]{\linewidth}\raggedright
\textbf{Qué puede adaptar SICHAR}
\end{minipage} \\
\midrule\noalign{}
\endfirsthead
\bottomrule\noalign{}
\endlastfoot
\textbf{Grupo Decor}

\textbf{Colombia}

\emph{\textbf{grupodecor.com}} & \begin{minipage}[t]{\linewidth}\raggedright
\begin{itemize}
\item
  Portafolio de acabados con red comercial que conecta a clientes de varias ciudades.
\item
  Atención diferenciada para distribuidores y para constructoras/proyectos.
\end{itemize}
\end{minipage} & Catálogo actualizado, \textbf{fotografías de proyectos reales, contenido sobre
combinaciones de acabados y campañas segmentadas} en redes, que generen oportunidades y faciliten
contactar a SICHAR o a sus distribuidores. \\
\textbf{Grainger}

\textbf{Internacional}

\emph{\textbf{grainger.com}} & \begin{minipage}[t]{\linewidth}\raggedright
\begin{itemize}
\item
  Facilita la compra B2B con catálogo y servicios digitales, sin depender solo de la atención
  presencial.
\end{itemize}
\end{minipage} & Contenido útil y de consulta (medidas, fichas técnicas, disponibilidad) que se
pueda \textbf{compartir y buscar}, además de publicaciones de inspiración. \\
\textbf{Cemex}

\textbf{México / Internacional}

\emph{\textbf{cemex.com}} & \begin{minipage}[t]{\linewidth}\raggedright
\begin{itemize}
\item
  Digitalizó por etapas: primero procesos internos, luego autoservicio y al final comunicación y
  redes.
\end{itemize}
\end{minipage} & Activar las redes \textbf{cuando ya haya procesos ordenados} detrás, para que cada
contacto reciba respuesta rápida. \\
\textbf{Corona / Houzz (referentes del hallazgo 3)}

\textbf{Colombia / EE. UU.} & \begin{minipage}[t]{\linewidth}\raggedright
\begin{itemize}
\item
  Usan proyectos reales e inspiración visual para que los usuarios descubran estilos y encuentren
  profesionales.
\end{itemize}
\end{minipage} & Mostrar \textbf{proyectos terminados} con productos SICHAR y dar crédito al
arquitecto o al distribuidor aliado. \\
\end{longtable}

\emph{Las empresas anteriores se toman como referentes de método digital. Antes de definir el
calendario final, se recomienda auditar la cuenta de Cerámica Italia (frecuencia, formatos, tipo de
contenido) como referente directo de competencia.}

\needspace{6\baselineskip}

\hypertarget{soluciuxf3n-tangible-para-sichar-5}{%
\subsection*{4. Solución tangible para SICHAR}\label{soluciuxf3n-tangible-para-sichar-5}}

Proponemos un \textbf{plan de contenido de 90 días} simple, sostenible con un equipo pequeño y
alineado con los perfiles del hallazgo 3.

\needspace{12\baselineskip}

\hypertarget{duxf3nde-estar-y-para-quiuxe9n}{%
\subsubsection*{Dónde estar y para quién}\label{duxf3nde-estar-y-para-quiuxe9n}}

\begin{longtable}[]{@{}
  >{\raggedright\arraybackslash}p{(\columnwidth - 4\tabcolsep) * \real{0.2183}}
  >{\raggedright\arraybackslash}p{(\columnwidth - 4\tabcolsep) * \real{0.3373}}
  >{\raggedright\arraybackslash}p{(\columnwidth - 4\tabcolsep) * \real{0.4444}}@{}}
\toprule\noalign{}
\begin{minipage}[b]{\linewidth}\raggedright
\textbf{Plataforma}
\end{minipage} & \begin{minipage}[b]{\linewidth}\raggedright
\textbf{Para quién}
\end{minipage} & \begin{minipage}[b]{\linewidth}\raggedright
\textbf{Qué se publica}
\end{minipage} \\
\midrule\noalign{}
\endfirsthead
\bottomrule\noalign{}
\endlastfoot
\textbf{Instagram} (principal) & Arquitectos, diseñadores, consumidor que remodela & Fotos y reels
de proyectos, combinaciones, producto insignia \\
\textbf{Pinterest} & Diseñadores y propietarios & Tableros por espacio y estilo (baños, cocinas,
fachadas) \\
\textbf{LinkedIn} & Constructoras y empresas de proyectos & Casos de obra, cumplimiento,
disponibilidad y alianzas \\
\textbf{WhatsApp Business} & Todos & Catálogo, atención rápida y derivación al distribuidor \\
\end{longtable}

\needspace{12\baselineskip}

\hypertarget{cuatro-pilares-de-contenido-y-calendario-semanal}{%
\subsubsection*{Cuatro pilares de contenido y calendario
semanal}\label{cuatro-pilares-de-contenido-y-calendario-semanal}}

\begin{longtable}[]{@{}
  >{\raggedright\arraybackslash}p{(\columnwidth - 6\tabcolsep) * \real{0.1290}}
  >{\raggedright\arraybackslash}p{(\columnwidth - 6\tabcolsep) * \real{0.2183}}
  >{\raggedright\arraybackslash}p{(\columnwidth - 6\tabcolsep) * \real{0.4345}}
  >{\raggedright\arraybackslash}p{(\columnwidth - 6\tabcolsep) * \real{0.2183}}@{}}
\toprule\noalign{}
\begin{minipage}[b]{\linewidth}\raggedright
\textbf{Día}
\end{minipage} & \begin{minipage}[b]{\linewidth}\raggedright
\textbf{Pilar}
\end{minipage} & \begin{minipage}[b]{\linewidth}\raggedright
\textbf{Ejemplo de publicación}
\end{minipage} & \begin{minipage}[b]{\linewidth}\raggedright
\textbf{Llamado a la acción}
\end{minipage} \\
\midrule\noalign{}
\endfirsthead
\bottomrule\noalign{}
\endlastfoot
\textbf{Lunes} & \textbf{Proyectos reales} & Antes/después de un baño o fachada con producto SICHAR
y crédito al distribuidor o arquitecto & ``Cotiza tu proyecto'' \\
\textbf{Miércoles} & \textbf{Guía de acabados} & Cómo combinar porcelanato brillante 60x130 con
grifería en tonos cálidos o negros & ``Descarga el catálogo'' \\
\textbf{Viernes} & \textbf{Producto insignia y novedades} & Video corto del porcelanato brillante
60x130 en distintos espacios & ``Dónde comprar'' \\
\textbf{Quincenal} & \textbf{Confianza y distribuidores} & Detrás de cámaras: llegada de
contenedores, control de calidad, presentación de un distribuidor aliado & ``Conoce a nuestros
aliados'' \\
\end{longtable}

\needspace{12\baselineskip}

\hypertarget{plan-de-acciuxf3n-3}{%
\subsubsection*{Plan de acción}\label{plan-de-acciuxf3n-3}}

\begin{longtable}[]{@{}
  >{\raggedright\arraybackslash}p{(\columnwidth - 8\tabcolsep) * \real{0.0437}}
  >{\raggedright\arraybackslash}p{(\columnwidth - 8\tabcolsep) * \real{0.4405}}
  >{\raggedright\arraybackslash}p{(\columnwidth - 8\tabcolsep) * \real{0.1637}}
  >{\raggedright\arraybackslash}p{(\columnwidth - 8\tabcolsep) * \real{0.1042}}
  >{\raggedright\arraybackslash}p{(\columnwidth - 8\tabcolsep) * \real{0.2480}}@{}}
\toprule\noalign{}
\begin{minipage}[b]{\linewidth}\raggedright
\textbf{\#}
\end{minipage} & \begin{minipage}[b]{\linewidth}\raggedright
\textbf{Acción concreta}
\end{minipage} & \begin{minipage}[b]{\linewidth}\raggedright
\textbf{Responsable}
\end{minipage} & \begin{minipage}[b]{\linewidth}\raggedright
\textbf{Plazo}
\end{minipage} & \begin{minipage}[b]{\linewidth}\raggedright
\textbf{Entregable}
\end{minipage} \\
\midrule\noalign{}
\endfirsthead
\bottomrule\noalign{}
\endlastfoot
1 & Auditar la cuenta actual de SICHAR y la de \textbf{Cerámica Italia}: frecuencia, formatos,
mensajes, interacción. & Líder digital & Semanas 1--2 & Informe comparativo \\
2 & Definir \textbf{identidad visual y tono de voz} (profesional, cercano, inspirador) y optimizar
el perfil: logo, biografía con propuesta de valor, enlace al catálogo y al localizador. & Líder
digital + diseño & Semanas 2--4 & Perfil optimizado y guía de estilo \\
3 & Armar un \textbf{banco de 30 fotografías y 10 videos cortos} de proyectos y productos (con
permiso de los clientes). & Comercial + apoyo creativo & Semanas 3--8 & Banco de contenido \\
4 & Publicar \textbf{3 veces por semana} según el calendario y responder mensajes en menos de 24
horas. & Líder digital & Desde el mes 3 & Calendario cumplido \\
5 & Probar una \textbf{pauta pequeña y segmentada} (arquitectos y constructoras en las ciudades de
las 5 sedes), con un presupuesto limitado a definir por la gerencia. & Líder digital + gerencia &
Meses 4--6 & Resultados de la pauta \\
6 & Revisar métricas cada mes y ajustar contenido; \textbf{eliminar lo que no funcione}. & Líder
digital & Mensual & Informe mensual \\
\end{longtable}

\needspace{12\baselineskip}

\hypertarget{cuxf3mo-medir-si-estuxe1-funcionando-5}{%
\subsection*{5. Cómo medir si está funcionando}\label{cuxf3mo-medir-si-estuxe1-funcionando-5}}

\begin{longtable}[]{@{}
  >{\raggedright\arraybackslash}p{(\columnwidth - 4\tabcolsep) * \real{0.3175}}
  >{\raggedright\arraybackslash}p{(\columnwidth - 4\tabcolsep) * \real{0.3452}}
  >{\raggedright\arraybackslash}p{(\columnwidth - 4\tabcolsep) * \real{0.3373}}@{}}
\toprule\noalign{}
\begin{minipage}[b]{\linewidth}\raggedright
\textbf{Indicador}
\end{minipage} & \begin{minipage}[b]{\linewidth}\raggedright
\textbf{Cómo se mide}
\end{minipage} & \begin{minipage}[b]{\linewidth}\raggedright
\textbf{Meta sugerida (a validar con la línea base)}
\end{minipage} \\
\midrule\noalign{}
\endfirsthead
\bottomrule\noalign{}
\endlastfoot
\textbf{Alcance y seguidores calificados} & Estadísticas de la plataforma; foco en arquitectos y
constructoras & Crecimiento sostenido mensual desde el mes 3 \\
\textbf{Mensajes y cotizaciones por redes} & Mensajes recibidos que piden precio o cotización &
Línea base en el mes 3; aumento trimestral \\
\textbf{Visitas a ``Dónde comprar'' / catálogo} & Clics desde redes hacia el sitio o WhatsApp &
Creciente desde el lanzamiento de la Vitrina \\
\textbf{Leads por segmento} & Contactos nuevos clasificados por perfil (hallazgo 3) & Primeros leads
de arquitectos y constructoras en 90 días \\
\end{longtable}

\begin{cajaalerta}{Precauciones y errores que debe evitar}

\begin{itemize}
\item
  \textbf{No publicar sin un destino claro}: toda pieza debe llevar a cotizar, descargar o ubicar un
  distribuidor.
\item
  Evitar mostrar precios finales al público que puedan entrar en conflicto con los distribuidores
  (ver hallazgo 1).
\item
  Pedir autorización a clientes y distribuidores antes de mostrar sus proyectos.
\end{itemize}

\end{cajaalerta}

\hypertarget{hoja-de-ruta-integrada}{%
\section{7. Hoja de ruta integrada}\label{hoja-de-ruta-integrada}}

Las seis soluciones se ejecutan \textbf{en secuencia lógica y en paralelo cuando es posible}. Esta
tabla resume qué se hace en cada horizonte para que la gerencia pueda dar seguimiento en una sola
página.

\begin{longtable}[]{@{}
  >{\raggedright\arraybackslash}p{(\columnwidth - 4\tabcolsep) * \real{0.2183}}
  >{\raggedright\arraybackslash}p{(\columnwidth - 4\tabcolsep) * \real{0.6429}}
  >{\raggedright\arraybackslash}p{(\columnwidth - 4\tabcolsep) * \real{0.1389}}@{}}
\toprule\noalign{}
\begin{minipage}[b]{\linewidth}\raggedright
\textbf{Horizonte}
\end{minipage} & \begin{minipage}[b]{\linewidth}\raggedright
\textbf{Qué se hace}
\end{minipage} & \begin{minipage}[b]{\linewidth}\raggedright
\textbf{Hallazgos}
\end{minipage} \\
\midrule\noalign{}
\endfirsthead
\bottomrule\noalign{}
\endlastfoot
\textbf{Primeros 30 días}

Diagnosticar y alinear & \begin{minipage}[t]{\linewidth}\raggedright
\begin{itemize}
\item
  Mapear el canal y entrevistar a 8--10 distribuidores clave.
\item
  Construir la hoja de costo puesto en bodega por referencia.
\item
  Designar al líder digital.
\item
  Agendar 15--20 entrevistas de buyer persona.
\end{itemize}
\end{minipage} & 1, 2, 3, 4 \\
\textbf{Días 31 a 90}

Decidir y construir bases & \begin{minipage}[t]{\linewidth}\raggedright
\begin{itemize}
\item
  Firmar y comunicar la política de canal y el programa Distribuidor Aliado.
\item
  Cotizar proveedores alternativos y renegociar con India.
\item
  Definir y aplicar la política de precios por escalones.
\item
  Redactar las fichas de buyer persona.
\item
  Crear el archivo maestro de productos, el catálogo digital y el CRM básico.
\item
  Iniciar el plan de contenido en redes (desde el mes 3).
\end{itemize}
\end{minipage} & 1, 2, 3, 4, 6 \\
\textbf{Meses 4 a 6}

Abrir la puerta & \begin{minipage}[t]{\linewidth}\raggedright
\begin{itemize}
\item
  Lanzar la Vitrina SICHAR (sitio, catálogo, formulario y localizador ``Dónde comprar'').
\item
  Lanzar el piloto de la Mesa de Asesoría en Cali.
\item
  Probar pauta pequeña y segmentada en redes.
\item
  Evaluar con el semáforo si se abre una Línea Directa limitada.
\end{itemize}
\end{minipage} & 1, 4, 5, 6 \\
\textbf{Meses 7 a 12}

Escalar lo que funciona & \begin{minipage}[t]{\linewidth}\raggedright
\begin{itemize}
\item
  Replicar la asesoría en las otras 4 sedes si el piloto lo justifica.
\item
  Activar el autoservicio para distribuidores y constructoras (Etapa 3).
\item
  Consolidar el tablero de indicadores por sede.
\item
  Revisar y actualizar la política de precios y de proveedores.
\end{itemize}
\end{minipage} & 2, 4, 5 \\
\end{longtable}

\needspace{6\baselineskip}

\hypertarget{gobierno-del-plan-quiuxe9n-da-seguimiento}{%
\subsubsection{Gobierno del plan: quién da
seguimiento}\label{gobierno-del-plan-quiuxe9n-da-seguimiento}}

\begin{itemize}
\item
  \textbf{Comité de seguimiento mensual} (gerencia general, dirección comercial, compras/finanzas y
  líder digital): revisa indicadores, desbloquea decisiones y ajusta metas.
\item
  \textbf{Tablero único de indicadores} con los KPI de las seis secciones, actualizado cada mes.
\item
  \textbf{Revisión trimestral} del semáforo de canal para decidir si se abre o se amplía la venta
  directa.
\end{itemize}

\hypertarget{pruxf3ximos-pasos-y-decisiones-que-necesitamos-de-usted}{%
\section{8. Próximos pasos y decisiones que necesitamos de
usted}\label{pruxf3ximos-pasos-y-decisiones-que-necesitamos-de-usted}}

Para arrancar en los próximos \textbf{30 días} necesitamos cinco decisiones de la gerencia:

\begin{enumerate}
\def\labelenumi{\arabic{enumi}.}
\item
  \textbf{Congelar el lanzamiento del B2C} hasta tener la política de canal firmada y comunicada a
  los distribuidores.
\item
  \textbf{Nombrar al líder digital} (puede ser de medio tiempo) y a un responsable de importaciones
  para la hoja de costos.
\item
  \textbf{Autorizar las entrevistas} a 8--10 distribuidores y 15--20 clientes y prospectos, con los
  coordinadores de sede.
\item
  \textbf{Definir un presupuesto inicial} para catálogo, sitio web sencillo, CRM y piloto de
  asesoría (a cotizar en las primeras semanas).
\item
  \textbf{Aprobar el comité de seguimiento mensual} y su calendario.
\end{enumerate}

\begin{cajamarca}{Mensaje final}

SICHAR no necesita dejar de ser quien es para crecer. Su mayor activo es la \textbf{confianza} que
ha construido durante más de 20 años con distribuidores y ferreterías. La estrategia propuesta crece
\textbf{con} ellos y no \textbf{contra} ellos: más visibilidad, más margen protegido, más servicio y
nuevos clientes, sin romper la promesa de \textbf{``Somos su proveedor, no su competidor''}.

\end{cajamarca}

\hypertarget{anexo-nota-metodoluxf3gica-y-fuentes}{%
\section{Anexo · Nota metodológica y fuentes}\label{anexo-nota-metodoluxf3gica-y-fuentes}}

\textbf{Origen de la información.} El diagnóstico se basa en el trabajo del equipo del Proyecto EA
con la empresa SICHAR (hallazgos positivos y negativos, análisis FODA) y en el documento de
\textbf{modelos inspiracionales} elaborado para cada hallazgo. La estructura del informe sigue la
metodología del reporte de análisis de Berry Healthy (2024) usado como ejemplo.

\textbf{Cifras de terceros.} Los precios y cifras de empresas de referencia son
\textbf{referenciales}, tomados de fuentes públicas de cada compañía, y deben verificarse antes de
usarse en decisiones comerciales. El ejemplo numérico del hallazgo 2 es \textbf{hipotético} y sirve
solo para explicar la mecánica del margen.

\textbf{Metas e indicadores.} Las metas son sugeridas y deben calibrarse con la línea base de SICHAR
durante el primer mes.

\needspace{12\baselineskip}

\hypertarget{sitios-consultados-para-los-modelos-inspiracionales}{%
\subsubsection{Sitios consultados para los modelos
inspiracionales}\label{sitios-consultados-para-los-modelos-inspiracionales}}

\begin{longtable}[]{@{}
  >{\raggedright\arraybackslash}p{(\columnwidth - 2\tabcolsep) * \real{0.2778}}
  >{\raggedright\arraybackslash}p{(\columnwidth - 2\tabcolsep) * \real{0.7222}}@{}}
\toprule\noalign{}
\begin{minipage}[b]{\linewidth}\raggedright
\textbf{Hallazgo}
\end{minipage} & \begin{minipage}[b]{\linewidth}\raggedright
\textbf{Empresas y fuentes}
\end{minipage} \\
\midrule\noalign{}
\endfirsthead
\bottomrule\noalign{}
\endlastfoot
\textbf{1. Conflicto de canal} & coartesas.com · listo.co · rocagroup.com · cat.com ·
porcelanosa.com · aqualisa.co.uk \\
\textbf{2. Costos de importación} & samisa.co · ikea.com · adidas-group.com · flooranddecor.com ·
Gano Excel (fuente no registrada) \\
\textbf{3. Buyer persona} & houzz.com · blog.hubspot.com · homedepot.com · procore.com · corona.co ·
easyline.com.co \\
\textbf{4. Atención especializada} & corona.co · corporate.homedepot.com · tileshop.com ·
porcelanosa.com · Grupo Bios (fuente no registrada) \\
\textbf{5 y 6. Digital y redes} & grainger.com · grupodecor.com · homecenter.com.co · cemex.com ·
comproagro.com \\
\end{longtable}
\end{document}
